\chapter{Annexe expérimentale}
\label{sec:supplement}

Pour chaque chapitre, on trouvera ici ses principales affirmations et ce sur quoi elles reposent: l'expérience où elles ont été mesurées, ce qui a été mesuré exactement et où cela a été publié. Le statut, dans la colonne de droite, indique la solidité de l'affirmation: \emph{mesuré}, il existe une expérience directe; \emph{théorie}, conséquence mathématique établie dans le chapitre ou dans un travail classique; \emph{modèle}, résultat d'un calcul sur le modèle du livre (scripts dans le dossier \texttt{models/}); \emph{estimation}, ordre de grandeur sans mesure directe; \emph{hypothèse}, explication plausible que l'expérience n'a pas encore démontrée. Là où il n'existe pas d'expérience, nous le disons.

\section{Chapitre~\ref{sec:neurontypes}. Types de neurones}

Les nombres du chapitre reposent sur des comptages directs de cellules dans le cerveau humain, là où ils existent; la règle \og un neurone, un neurotransmetteur\fg, sur des atlas cellulaires et sur des expériences qui en ont aussi trouvé les exceptions; le rôle de \nt{DOPA}, sur des enregistrements d'impulsions; et les rôles des autres canaux de diffusion, sur des hypothèses.

{\footnotesize
\setlength{\tabcolsep}{3pt}\renewcommand{\arraystretch}{1.15}
\begin{longtable}{>{\raggedright}p{0.5cm}>{\raggedright}p{4.0cm}>{\raggedright}p{5.6cm}>{\raggedright}p{2.3cm}>{\raggedright\arraybackslash}p{1.7cm}}
\toprule
N\textsuperscript{o} & Affirmation & Expérience et ce qui est mesuré & Source & Statut \\
\midrule\endfirsthead
\toprule
N\textsuperscript{o} & Affirmation & Expérience et ce qui est mesuré & Source & Statut \\
\midrule\endhead
1 & Le cerveau humain compte environ $8.6\cdot10^{10}$ neurones, et presque tous les neurones \nt{GLUT} sont les petits neurones du cervelet (tabl.~\ref{tab:types}) & Fractionneur isotrope: le tissu est dissous en une suspension homogène de noyaux, et l'on compte les noyaux porteurs du marqueur neuronal NeuN. Chez l'homme adulte, $86.1\pm8.1$ milliards de neurones; le cortex cérébral n'en contient que $19\,\%$, la masse principale est dans le cervelet & \cite{Azevedo2009} & mesuré \\
2 & Presque chaque neurone a un seul neurotransmetteur principal (proposition~\ref{prop:dale}), mais il y a des exceptions & Atlas transcriptomique du cerveau de la souris: environ $4\cdot10^6$ cellules, $5322$ groupes; chaque groupe reçoit un neurotransmetteur d'après les gènes de synthèse et d'empaquetage. Les exceptions sont mesurées directement: en tranches, les neurones \nt{DOPA} libèrent aussi du GABA par le même transporteur vésiculaire et inhibent les neurones du striatum; dans une partie des axones \nt{DOPA}, le glutamate et la dopamine sortent de terminaisons différentes d'un même fil (microscopie électronique et optogénétique) & \cite{Strata1999,Yao2023,Tritsch2012,Zhang2015} & mesuré \\
3 & Toute la colonne de la matrice des poids est de même signe~\eqref{eq:dalesign} & Déduction du chapitre à partir de la proposition~\ref{prop:dale} et du fait que presque tous les récepteurs du glutamate sont excitateurs et ceux du GABA inhibiteurs. Il n'existe pas de vérification directe de \og tous les destinataires d'un neurone reçoivent un poids de même signe\fg{} comme règle générale; contre-exemples: la co-libération de GABA par les neurones \nt{DOPA} (ligne~2) et les destinataires dotés de récepteurs de signes différents (ligne~11) & déduction du chapitre & théorie \\
4 & De trois quarts à quatre cinquièmes des neurones du cortex sont \nt{GLUT}, d'un cinquième à un quart sont \nt{GABA} & Immunomarquage du GABA dans des colonnes de $50$\,\textmu m de large traversant toute l'épaisseur du cortex, dans $10$ aires corticales du singe: dans $8$ aires, les neurones GABA forment environ $25\,\%$ de tous les neurones, dans le cortex visuel primaire moins de $20\,\%$ & \cite{Hendry1987} & mesuré \\
5 & Un neurone \nt{GLUT} a environ $10^4$ sorties & Stéréologie post mortem: dans le néocortex de jeunes hommes, $1.64\cdot10^{14}$ synapses (microscopie électronique, dissecteur) et environ $2.3\cdot10^{10}$ neurones. Le rapport donne $\sim7\cdot10^3$ synapses par neurone; en moyenne, le nombre d'entrées égale le nombre de sorties & \cite{Tang2001synapses,Pakkenberg1997} & mesuré \\
6 & La sortie d'un neurone \nt{DOPA} se ramifie en $10^5$--$10^6$ terminaisons; c'est une estimation d'après la couverture de la cible, pas un comptage & Marquage viral de la membrane de neurones \nt{DOPA} isolés du rat et reconstruction complète de l'axone: un neurone couvre de $0.45$ à $5.7\,\%$ (en moyenne $2.7\,\%$) du volume du striatum. C'est la couverture qui est mesurée, pas le nombre de terminaisons: celui-ci en est déduit en ordre de grandeur, sans comptage direct. Pour \nt{SERO} et \nt{NORA}, les nombres de terminaisons sont des estimations grossières & \cite{Matsuda2009} & mesuré \\
7 & Les neurones de diffusion sont peu nombreux: \nt{DOPA} $\sim5\cdot10^5$ (le principal des deux groupes, borne inférieure), \nt{SERO} $\sim3\cdot10^5$ (somme de deux comptages partiels), \nt{HIST} $\sim6\cdot10^4$ (tabl.~\ref{tab:types}, fig.~\ref{fig:typepie}) & Comptages stéréologiques chez l'homme: $550\,000$ neurones pigmentés dans la substance noire (sans l'aire tegmentale ventrale); $235\,000$ neurones dans le noyau dorsal du raphé (tous les neurones du noyau) et encore $125\,000$ neurones sérotoninergiques dans le tegmentum du pont; environ $64\,000$ neurones histaminergiques de l'hypothalamus. Pour \nt{NORA}, \nt{GLYC}, \nt{ACET} et \nt{ADRE}, pas de comptage direct: le tableau donne des estimations grossières en ordre de grandeur & \cite{Pakkenberg1991,Baker1990,Baker1991,Airaksinen1991} & mesuré \\
8 & Le glutamate dans la fente monte jusqu'à environ une millimole par litre et retombe en $\sim1\ms$ & D'après le déplacement d'un bloqueur compétitif à dissociation rapide des récepteurs NMDA pendant la transmission synaptique (culture de neurones d'hippocampe): pic de $1.1$\,mM, constante de décroissance de $1.2\ms$ & \cite{Clements1992} & mesuré \\
9 & Dans le cortex en activité, les courants excitateur et inhibiteur sont presque équilibrés, et le neurone reste près du seuil & Théorie: un réseau à connexions fortes atteint de lui-même le régime équilibré. Expérience: dans le cortex préfrontal du furet in vivo, pendant les états \og up\fg, le potentiel d'inversion du courant synaptique reste vers $-37$\,mV pendant des centaines de millisecondes, alors que la conductance varie de $\sim21$\,nS: excitation et inhibition montent et descendent ensemble & \cite{vanVreeswijk1996,Haider2006} & mesuré \\
10 & Les neurones \nt{DOPA} annoncent l'erreur de prédiction de la récompense~\eqref{eq:rpe} (fig.~\ref{fig:rpe}) & Impulsions des neurones \nt{DOPA} du singe: bouffée au jus inattendu; après apprentissage, bouffée au signal et aucune réponse au jus prédit; creux de fréquence quand le jus promis n'est pas donné. Dans l'expérience quantitative, la fréquence est proportionnelle à la différence entre la récompense actuelle et la moyenne pondérée des précédentes, mais seulement pour les erreurs positives. L'équation~\eqref{eq:rpe} elle-même est l'algorithme des différences temporelles & \cite{Schultz1997,Bayer2005,SuttonBarto2018} & mesuré \\
11 & Le signe et la vitesse sont fixés par le récepteur: ionotropes, fractions de milliseconde; métabotropes, beaucoup plus lents (proposition~\ref{prop:receptor}) & Brèves impulsions de glutamate sur des fragments de membrane de neurones d'hippocampe: le courant des récepteurs ionotropes monte (de $20$ à $80\,\%$) en $0.2$--$0.6\ms$ et retombe en $2$--$3\ms$. Le potentiel inhibiteur lent des neurones pyramidaux est supprimé par un bloqueur sélectif du récepteur métabotrope du GABA. Deux familles de récepteurs de la dopamine à action opposée; dans le striatum, les neurones à récepteur D1 ont besoin de dopamine pour renforcer les synapses, ceux à D2 pour les affaiblir & \cite{Colquhoun1992,DutarNicoll1988,Beaulieu2011,Shen2008} & mesuré \\
12 & Un canal de diffusion n'est pas un fil unique, mais un faisceau de quelques lignes (section~\ref{sec:buses}) & Marquage viro-génétique des neurones sérotoninergiques du noyau dorsal du raphé de la souris: les neurones qui projettent vers l'amygdale et vers le cortex frontal occupent des places différentes du noyau, se ramifient vers des régions complémentaires et répondent de façon opposée à un stimulus désagréable. Revue: des sous-groupes de neurones du locus coeruleus (\nt{NORA}) codent des processus différents & \cite{Ren2018,Poe2020} & mesuré \\
13 & \nt{NORA} fixe le gain $g$ & Hypothèse du gain adaptatif; modèle de réseau dans lequel les catécholamines augmentent la pente de la caractéristique de transfert. Pas de mesure directe de \og $g$ croît avec la noradrénaline\fg{} à l'échelle du réseau entier; mesuré: le diamètre de la pupille suit les impulsions des neurones du locus coeruleus du singe & \cite{AstonJones2005,ServanSchreiber1990,Joshi2016} & hypothèse \\
14 & \nt{ACET} commute \og écriture/lecture\fg{} et fixe la vitesse d'apprentissage & Hypothèse. Appui: dans des tranches de cortex olfactif du rat, les agonistes de l'acétylcholine ($100$\,\textmu M) affaiblissent les synapses internes \og de rappel\fg{} de plus de $60\,\%$, et les synapses d'entrée de moins de $15\,\%$ & \cite{Hasselmo2006,HasselmoBower1992} & hypothèse \\
15 & \nt{SERO} fixe le facteur d'actualisation $\gamma$; les neurones sérotoninergiques travaillent régulièrement, à quelques impulsions par seconde, et se taisent presque pendant l'une des phases du sommeil & Hypothèse \og sérotonine $=\gamma$\fg. Argument le plus fort: l'activation optogénétique des neurones sérotoninergiques du noyau dorsal du raphé de la souris réduit le nombre d'abandons prématurés de l'attente d'une récompense différée. La fréquence des impulsions et le silence pendant le sommeil à mouvements oculaires rapides sont mesurés (revue d'enregistrements) & \cite{Doya2002,Miyazaki2014,JacobsAzmitia1992} & hypothèse \\
\bottomrule
\end{longtable}}

\section{Chapitre~\ref{sec:glutcomputer}. Ce que calculent les neurones \nt{GLUT}}

Les énoncés logiques du chapitre sont des théorèmes sur les circuits à poids positifs ou nuls, démontrés dans le chapitre lui-même; l'expérience confirme le modèle de neurone et le fait que les circuits construits d'après ces théorèmes (détecteur de coïncidences, ligne à retard, groupe auto-entretenu, chaîne) se rencontrent réellement dans le cerveau.

{\footnotesize
\setlength{\tabcolsep}{3pt}\renewcommand{\arraystretch}{1.15}
\begin{longtable}{>{\raggedright}p{0.5cm}>{\raggedright}p{4.0cm}>{\raggedright}p{5.6cm}>{\raggedright}p{2.3cm}>{\raggedright\arraybackslash}p{1.7cm}}
\toprule
N\textsuperscript{o} & Affirmation & Expérience et ce qui est mesuré & Source & Statut \\
\midrule\endfirsthead
\toprule
N\textsuperscript{o} & Affirmation & Expérience et ce qui est mesuré & Source & Statut \\
\midrule\endhead
1 & Le neurone est un circuit RC avec seuil et remise à zéro~\eqref{eq:glutneuron} & On a injecté dans le corps de neurones pyramidaux du cortex du rat (tranches) un courant bruité semblable au flux synaptique du cerveau vivant. La dépendance de la fréquence moyenne à la moyenne et à la dispersion du courant est décrite par un neurone intégrateur à fuite, si l'on ajoute une lente adaptation de la fréquence. Les plages de $\tau$ et des retards sont données dans le chapitre sans référence expérimentale & \cite{Rauch2003} & mesuré \\
2 & Le modèle~\eqref{eq:glutneuron} est une simplification: les branches d'entrée émettent elles-mêmes des impulsions locales & Enregistrement directement sur les dendrites de neurones pyramidaux du cortex visuel de la souris in vivo: un stimulus visuel déclenche des impulsions dendritiques locales dépendant des récepteurs NMDA; les supprimer élargit l'accord du neurone sur l'orientation & \cite{Smith2013} & mesuré \\
3 & Fenêtre de coïncidence $\Delta_{\max}=\tau\ln\frac{w}{\theta-w}$~\eqref{eq:window}; pour $\theta=1.5w$, elle vaut $0.7\tau$ & Conséquence du modèle~\eqref{eq:glutneuron}. Une mesure directe d'une fenêtre de sommation étroite existe pour les neurones à inhibition rapide (chapitre~\ref{sec:gaba}, ligne~3 de son tableau) & déduction du chapitre & théorie \\
4 & Un réseau fait seulement de neurones \nt{GLUT} ne calcule que des fonctions monotones des entrées (proposition~\ref{prop:monotone}) & Démonstration dans le chapitre: itération à partir du silence avec un membre de droite non décroissant. C'est un énoncé sur le modèle; aucune expérience ne l'a vérifié sur un réseau vivant privé d'inhibition & déduction du chapitre & théorie \\
5 & Les organes des sens fournissent une paire de fils: certaines cellules répondent au renforcement du stimulus, d'autres à son affaiblissement & Rétine: dès les cellules bipolaires, le signal des cônes se sépare en canaux ON et OFF. Blocage pharmacologique des cellules bipolaires ON chez le singe: les canaux restent séparés jusqu'au cortex; après le blocage, la détection d'une augmentation de lumière se dégrade, mais pas celle d'une diminution & \cite{Schiller1992} & mesuré \\
6 & Avec le code à deux fils, un réseau de neurones \nt{GLUT} calcule n'importe quelle fonction logique de façon asynchrone, sans cadence commune, au prix de deux fois plus de neurones (proposition~\ref{prop:dualrail}, \eqref{eq:dualrail}, fig.~\ref{fig:dualxor}) & Le code à deux fils et la détection de fin de calcul par le signal lui-même sont une technique standard des circuits asynchrones insensibles aux retards. Le circuit OU exclusif à six neurones est construit dans le chapitre. Aucune expérience ne montre que le cerveau calcule les fonctions logiques précisément en code à deux fils & \cite{Sparso2001} & théorie \\
7 & La plus grande différence de temps d'arrivée du son aux deux oreilles est chez l'homme de $\approx0.6\ms$, et le cerveau distingue environ dix microsecondes & Auditeurs dans une expérience à choix forcé entre deux réponses: seuil de discrimination de la différence de temps (75\,\% de réponses correctes) de $9$\,\textmu s pour un bruit dans la bande $150$--$1700$\,Hz, $11$\,\textmu s pour un son pur de $1$\,kHz, $28$\,\textmu s pour un clic. La limite de $0.6\ms$ est un calcul du chapitre, $0.22\,\text{m}/343\,\text{m/s}$ & \cite{KlumppEady1956} & mesuré \\
8 & Chez les oiseaux, la différence de temps se change en position grâce à des lignes à retard et des détecteurs de coïncidences~\eqref{eq:itd} (fig.~\ref{fig:itd}) & Effraie des clochers: les axones venus des deux oreilles entrent dans le noyau des détecteurs de coïncidences par des côtés opposés; le temps d'arrivée des impulsions le long de l'axone varie de façon ordonnée avec la profondeur, et l'étendue des retards à travers le noyau ($160$\,\textmu s sur $700$\,\textmu m) correspond à la plage des retards interauraux & \cite{Carr1990} & mesuré \\
9 & Chez les mammifères, on n'a pas trouvé de rangée de retards; la même tâche est résolue par des détecteurs de coïncidences avec une inhibition précisément calée venue de neurones \nt{GLYC} & Enregistrement in vivo dans l'olive supérieure médiane de la gerbille: les réponses ne forment pas de carte des directions; l'application de glycine et de son bloqueur par micropipette montre que l'inhibition \emph{glycinergique} précisément calée dans le temps fixe la plage utile des retards. L'inhibition vient ici de \nt{GLYC}, non de \nt{GABA} & \cite{Brand2002,Grothe2010} & mesuré \\
10 & Un groupe à forte rétroaction est une bascule; l'activité auto-entretenue persiste des secondes après le stimulus (fig.~\ref{fig:bistable}) & Cortex préfrontal du singe dans une tâche de saccade différée vers un point mémorisé (délai de $1$--$6$\,s): $87$ des $170$ neurones liés à la tâche changent de fréquence pendant le délai, et pour $79\,\%$ d'entre eux cela dépend de la direction du point mémorisé. Que cette activité soit entretenue par une rétroaction à deux états stables relève de la théorie & \cite{Funahashi1989,Wang2001} & mesuré \\
11 & Le bit s'écrit si l'apport dure plus de $\approx0.7\tau$ (fig.~\ref{fig:recurrentgroup}b) & Calcul de l'équation $\tau\,\dd r/\dd t=-r+f(wr+I)$ par la méthode d'Euler pour $w=1$, $I=0.6$; pas de script dans \texttt{models/}. Pas d'expérience & calcul du chapitre & modèle \\
12 & La mémoire peut être gardée dans la facilitation des synapses, presque sans impulsions & Théorie: le calcium résiduel dans les terminaisons garde l'information environ une seconde. Appui: dans le cortex préfrontal médian du furet (enregistrement de plusieurs neurones à la fois), il existe un sous-réseau de neurones pyramidaux reliés par des synapses facilitatrices à long effet rémanent. Revue des observations d'intervalles \og silencieux\fg{} pendant le maintien en mémoire. Quelle façon le cortex utilise et quand reste une question ouverte & \cite{Mongillo2008,WangMarkram2006,Stokes2015} & hypothèse \\
13 & La mémoire s'efface par fatigue: la réserve de neurotransmetteur s'épuise & Enregistrements appariés de neurones pyramidaux du cortex: la réponse de la synapse baisse lors d'impulsions fréquentes, la vitesse de la baisse est fixée par la probabilité de libération. Qu'elle éteigne précisément un groupe auto-entretenu n'est pas montré directement & \cite{Tsodyks1997} & mesuré \\
14 & Une chaîne de groupes est une minuterie déclenchée par un événement; chez les oiseaux chanteurs, les neurones se déclenchent strictement à tour de rôle & Enregistrement de neurones identifiés du centre vocal supérieur du diamant mandarin pendant le sommeil et le chant: chaque neurone qui projette vers le noyau moteur émet une seule brève bouffée à un instant précis du chant, et les neurones se déclenchent successivement & \cite{Hahnloser2002} & mesuré \\
\bottomrule
\end{longtable}}

\section{Chapitre~\ref{sec:gaba}. \nt{GLUT} et \nt{GABA} ensemble}

Les énoncés logiques et dynamiques du chapitre y sont déduits de l'analyse linéaire et de la loi d'Ohm; l'expérience montre que les circuits sur lesquels ils reposent (veto retardé, inhibition directe à la suite de l'excitation, stabilisation par l'inhibition, rythme de la boucle \nt{GLUT}--\nt{GABA}) fonctionnent réellement dans le cerveau.

{\footnotesize
\setlength{\tabcolsep}{3pt}\renewcommand{\arraystretch}{1.15}
\begin{longtable}{>{\raggedright}p{0.5cm}>{\raggedright}p{4.0cm}>{\raggedright}p{5.6cm}>{\raggedright}p{2.3cm}>{\raggedright\arraybackslash}p{1.7cm}}
\toprule
N\textsuperscript{o} & Affirmation & Expérience et ce qui est mesuré & Source & Statut \\
\midrule\endfirsthead
\toprule
N\textsuperscript{o} & Affirmation & Expérience et ce qui est mesuré & Source & Statut \\
\midrule\endhead
1 & Les neurones \nt{GABA} représentent d'un cinquième à un quart des neurones du cortex & Immunomarquage du GABA dans $10$ aires du cortex du singe: environ $25\,\%$ des neurones dans $8$ aires, moins de $20\,\%$ dans le cortex visuel primaire. Revue de la diversité des neurones inhibiteurs du cortex & \cite{Hendry1987,Markram2004} & mesuré \\
2 & Ce que fait l'inhibition dépend en grande partie du site d'arrivée: dendrite, corps, lieu de naissance de l'impulsion, terminaison du fil d'un autre neurone & Revue: les classes de neurones \nt{GABA} du cortex se distinguent par leur cible sur le destinataire. Enregistrement simultané du corps et d'une dendrite d'un neurone pyramidal: l'inhibition directe est bien plus forte sur le corps que sur les dendrites. L'inhibition sur la terminaison du fil d'un autre neurone, dans la moelle épinière, réduit la libération de neurotransmetteur d'une seule ligne d'entrée (revue de mesures) & \cite{Markram2004,Pouille2001,Rudomin1999} & mesuré \\
3 & Le changement de signe par un intermédiaire coûte un retard, environ une milliseconde; l'inhibition qui suit l'excitation rétrécit la fenêtre de coïncidence & Neurones pyramidaux de l'hippocampe du rat en tranches: l'inhibition par un intermédiaire suit l'excitation directe avec un court retard, et la fenêtre dans laquelle deux entrées excitatrices se somment jusqu'au seuil est inférieure à $2\ms$. Sur les dendrites, où cette inhibition est plus faible, la fenêtre est plus large & \cite{Pouille2001} & mesuré \\
4 & Le veto $a\wedge\bar b$~\eqref{eq:veto}, avec le OU et une constante égale à un, est fonctionnellement complet (proposition~\ref{prop:gabacomplete}) & Démonstration dans le chapitre. Résultat classique: un réseau d'éléments à seuil dotés d'entrées inhibitrices réalise n'importe quelle expression logique. Énoncé sur le modèle, pas d'expérience & \cite{McCullochPitts1943} & théorie \\
5 & Le veto retardé donne un détecteur de direction du mouvement de vitesse optimale $v^*=\Delta x/d$~\eqref{eq:vstar} & Rétine du lapin: la sélectivité à la direction est créée par une inhibition qui, pour un mouvement dans la direction \og nulle\fg, éteint l'excitation. L'enregistrement séparé des entrées excitatrice et inhibitrice des cellules sélectives a montré que les cellules amacrines en étoile (\nt{GABA}) les inhibent de façon asymétrique, seulement par les prolongements orientés du côté \og nul\fg. La formule de $v^*$ est une déduction du chapitre & \cite{Barlow1965,Fried2002} & mesuré \\
6 & L'inhibition shuntante divise la tension au lieu de soustraire~\eqref{eq:shunt} (fig.~\ref{fig:shunt}) & Loi d'Ohm pour un diviseur à trois conductances, avec un potentiel d'équilibre du chlorure proche du repos; pour un neurone sans impulsions & déduction du chapitre & théorie \\
7 & Sur la fréquence des impulsions, le shunt agit par soustraction, et la division vient de l'inhibition associée à un bruit de fond équilibré & Modèle: chez un neurone qui se déclenche, le courant à travers le shunt est presque constant, et l'effet sur la fréquence est soustractif. Expérience: dans des neurones pyramidaux du cortex somatosensoriel du rat (tranches), on a introduit par clamp dynamique des conductances de fond excitatrice et inhibitrice; leur augmentation simultanée et équilibrée divise la pente de la réponse sans décalage notable de la fréquence. Revue: la division par l'activité totale se rencontre dans beaucoup de régions du cerveau & \cite{HoltKoch1997,Chance2002,Carandini2012} & mesuré \\
8 & Pour $\beta>1$, l'inhibition mutuelle choisit un seul gagnant (proposition~\ref{prop:wta}, \eqref{eq:wta}) & Valeurs propres $-(1\pm\beta)/\tau$: déduction du chapitre. Les fréquences $0.73$ et $0.53$ pour $\beta=0.5$ sont le point fixe $r_{1,2}=(I_{1,2}-\beta I_{2,1})/(1-\beta^2)$; pas de script dans \texttt{models/}. Le chapitre ne cite pas d'expérience directe & déduction du chapitre & théorie \\
9 & Effacement de la mémoire par un signal via \nt{GABA}; deux groupes à inhibition mutuelle forment un verrou & Découle de la fig.~\ref{fig:bistable} et de la proposition~\ref{prop:wta}. Pas d'expérience & déduction du chapitre & théorie \\
10 & L'équilibre de la boucle \nt{GLUT}--\nt{GABA} est stable si l'inhibition est assez forte et assez rapide (proposition~\ref{prop:ei}) & Modèle à deux populations~\eqref{eq:ei}; les conditions (i) et (ii) portent sur le déterminant et la trace de sa matrice, déduits dans le chapitre & \cite{WilsonCowan1972} & théorie \\
11 & Pour $w_{EE}=1.5$ et $w_{EI}w_{IE}=2$, une inhibition rapide ($\tau_I=2\ms$) maintient $E^*=0.67h$, une inhibition lente ($30\ms$) déclenche des oscillations (fig.~\ref{fig:ei}) & Résolution numérique de~\eqref{eq:ei}, script \texttt{figures/make\_ei.py} & calcul & modèle \\
12 & Le cortex fonctionne en régime stabilisé par l'inhibition; signature: poussez les neurones \nt{GABA}, et leur fréquence baisse~\eqref{eq:isn} & La théorie a prédit la réponse paradoxale. Expérience: enregistrements intracellulaires dans le cortex visuel primaire du chat; quand la réponse est supprimée par un stimulus environnant, la conductance inhibitrice n'augmente pas, mais baisse avec la conductance excitatrice. L'excitation optogénétique des neurones PV de la souris dans les cortex visuel, somatosensoriel et moteur abaisse la fréquence des neurones inhibiteurs, y compris sans stimulus et sous anesthésie légère. Le coefficient $-1/3$ avec les valeurs de la fig.~\ref{fig:ei} est une déduction du chapitre & \cite{Tsodyks1997isn,Ozeki2009,Sanzeni2020} & mesuré \\
13 & La boucle \og \nt{GLUT} excite \nt{GABA}, \nt{GABA} inhibe \nt{GLUT}\fg{} engendre un rythme rapide de période $12$--$50\ms$ & La stimulation optogénétique des neurones \nt{GABA} rapides dans le cortex somatosensoriel de la souris, à des fréquences de $8$--$200$\,Hz, renforce sélectivement le rythme gamma ($20$--$80$\,Hz, période $12$--$50\ms$); la stimulation des neurones \nt{GLUT} ne renforce que des rythmes plus lents & \cite{Cardin2009,Buzsaki2012} & mesuré \\
\bottomrule
\end{longtable}}

\section{Chapitre~\ref{sec:code}. Le code Morse}

Les bornes du chapitre relèvent de la théorie de l'information et du calcul; sont mesurés l'endroit où naît le bruit dans le canal, la vitesse du cerveau et le coût d'une impulsion, tandis que la question du code que le cortex utilise réellement reste ouverte.

{\footnotesize
\setlength{\tabcolsep}{3pt}\renewcommand{\arraystretch}{1.15}
\begin{longtable}{>{\raggedright}p{0.5cm}>{\raggedright}p{4.0cm}>{\raggedright}p{5.6cm}>{\raggedright}p{2.3cm}>{\raggedright\arraybackslash}p{1.7cm}}
\toprule
N\textsuperscript{o} & Affirmation & Expérience et ce qui est mesuré & Source & Statut \\
\midrule\endfirsthead
\toprule
N\textsuperscript{o} & Affirmation & Expérience et ce qui est mesuré & Source & Statut \\
\midrule\endhead
1 & Les fréquences des neurones \nt{GLUT} du cortex vont d'une fraction à quelques dizaines d'impulsions par seconde, et la plupart du temps les neurones travaillent près de la borne inférieure & Enregistrement par une pipette accolée à la cellule (le choix du neurone ne dépend pas de son activité) dans le cortex auditif du rat éveillé: à tout instant, moins de $5\,\%$ des neurones dépassent $20$ impulsions par seconde; chez une partie des neurones, la fréquence de fond est inférieure à $0.01$ par seconde; la distribution des fréquences est log-normale & \cite{Hromadka2008} & mesuré \\
2 & Capacité d'un canal à impulsions $R_{\max}\approx r\log_2\frac{e}{r\Delta t}$~\eqref{eq:spikeentropy}, presque entièrement dans l'instant des impulsions (proposition~\ref{prop:capacity}) & Entropie d'une chaîne binaire de cases de longueur $\Delta t$. Valeurs du chapitre: $8$ bits par impulsion pour $r=10$ par seconde et $\Delta t=1\ms$, environ $5$ bits pour $\Delta t=10\ms$, moins de $2$ bits par seconde pour un compteur d'impulsions & \cite{MacKay1952} & théorie \\
3 & Les neurones sensoriels transmettent d'un à quelques bits par impulsion, dans les meilleurs cas jusqu'à la moitié de la limite & Reconstruction du stimulus à partir des impulsions de neurones sensoriels dont le stimulus est connu; synthèse des mesures dans le livre & \cite{Rieke1997} & mesuré \\
4 & Le générateur d'impulsions est précis; ce sont les variations rapides de l'entrée qui fixent l'instant précis & Neurones du cortex du rat en tranches: pour un courant répété à fluctuations rapides, les impulsions se répètent à mieux que $1\ms$ près; pour un courant constant, les instants se dispersent & \cite{Mainen1995} & mesuré \\
5 & La synapse n'est pas fiable: plus de la moitié des impulsions n'atteignent pas le destinataire, à cause du caractère aléatoire de la libération & Stimulation minimale des entrées de neurones pyramidaux de l'hippocampe (tranches): plus de la moitié des impulsions ne provoquent pas de réponse. Sont exclues les fluctuations du seuil de stimulation, les défauts de conduction aux ramifications de l'axone et la basse température de l'expérience; il reste la libération probabiliste & \cite{AllenStevens1994} & mesuré \\
6 & Le flux d'impulsions du cortex vivant paraît aléatoire: intervalles dispersés comme dans un flux de Poisson, variance du nombre d'impulsions proche de la moyenne; une partie du \og bruit\fg{} est un message sur les mouvements de l'animal & Enregistrements dans les aires visuelles V1 et MT du singe éveillé: coefficient de variation des intervalles proche de $1$, rapport de la variance du nombre d'impulsions à la moyenne comme pour un processus aléatoire; un modèle qui somme de petites entrées indépendantes donne une dispersion bien plus faible. Dans le cortex visuel de la souris, une part importante de la variabilité se prédit d'après l'enregistrement vidéo des mouvements de l'animal lui-même & \cite{Softky1993,Stringer2019} & mesuré \\
7 & Le code de fréquence est lent: erreur $1/\sqrt{rT}$~\eqref{eq:rateerror}, alors que le cerveau reconnaît une image en $\sim150\ms$ & La formule relève de la statistique de Poisson, déduction du chapitre. Expérience: des sujets décidaient si un animal figurait sur une photographie inconnue montrée pendant $20\ms$; les potentiels cérébraux pour \og oui\fg{} et \og non\fg{} divergent au bout d'environ $150\ms$. Le découpage de ce temps en une dizaine d'étapes de $10$--$20\ms$ est une estimation du chapitre, pas une mesure & \cite{Thorpe1996} & mesuré \\
8 & Le moyennage sur un groupe est limité par la corrélation: l'erreur ne baisse pas de plus d'un facteur $1/\sqrt c$~\eqref{eq:correlated} (proposition~\ref{prop:pooling}) & Paires de neurones de l'aire MT du singe enregistrés simultanément: coefficient de corrélation des nombres d'impulsions de $0.12$ en moyenne, et l'analyse théorique montre que même cette corrélation limite le gain du groupe. Théorie: seule la part de la corrélation qui ressemble au signal impose une limite. Modèle de la position opposée: la fréquence d'un groupe de $50$--$100$ neurones se lit en un seul intervalle entre impulsions, $10$--$50\ms$, et l'instant des impulsions isolées n'est pas utilisé dans le cortex & \cite{Zohary1994,MorenoBote2014,ShadlenNewsome1998} & mesuré \\
9 & Un nombre peut s'écrire dans l'instant de la première impulsion~\eqref{eq:latency}, dans l'ordre des impulsions d'un groupe ou dans la phase d'un rythme local & Rétine de la salamandre: la structure spatiale d'une image brièvement présentée est inscrite dans les délais relatifs des premières impulsions des neurones de sortie, presque indépendamment du contraste. Cellules de lieu de l'hippocampe du rat: en traversant \og leur\fg{} lieu, les impulsions se décalent de plus en plus tôt dans le cycle du rythme thêta ($7$--$12$\,Hz), avec un décalage de $100^\circ$ à $355^\circ$. Dans quelle mesure le cortex utilise l'instant des impulsions reste une question ouverte (ligne~8) & \cite{Gollisch2008,OKeefe1993} & mesuré \\
10 & Le code lu dépend de la constante de temps du destinataire~\eqref{eq:reader} (proposition~\ref{prop:reader}); dans le cortex actif, les neurones sont plus proches de détecteurs de coïncidences & La formule~\eqref{eq:reader} est une déduction du chapitre. Revue d'enregistrements in vivo: dans le cortex actif, la conductance de la membrane est plusieurs fois plus élevée qu'au repos, et la constante de temps d'autant plus courte. Que le cortex change de code précisément ainsi n'est pas montré directement & \cite{Destexhe2003} & hypothèse \\
11 & Une synapse à dépression transmet la variation de fréquence, au fond $\Delta r/r$~\eqref{eq:depression} (fig.~\ref{fig:depression}) & Enregistrements appariés de neurones pyramidaux du cortex: la vitesse de dépression est fixée par la probabilité de libération; avec une dépression lente, la réponse reflète la fréquence, avec une dépression rapide, les coïncidences. Modèle fondé sur la dépression mesurée dans le cortex du rat: des variations relatives égales de fréquence sur des entrées rapides et lentes donnent la même réponse, c'est-à-dire que la synapse transmet $\Delta r/r$. Les nombres $1.7\to2.2$, $6.7$ et $90\ms$ sont des calculs du chapitre & \cite{Tsodyks1997,Abbott1997depression} & mesuré \\
12 & La bouffée est un \og trait\fg: sa probabilité de se perdre, $(1-p)^k$~\eqref{eq:burst}, est faible, et la facilitation la réduit encore & Revue: aux synapses peu fiables, les impulsions isolées se perdent souvent, tandis que les bouffées sont transmises de façon fiable grâce à la facilitation; une bouffée provoque des modifications durables des synapses. Que le cerveau lise la bouffée comme un signe à part est une hypothèse & \cite{Lisman1997,AllenStevens1994} & hypothèse \\
13 & Les neurones \nt{GABA} rapides fixent le rythme et la \og ponctuation\fg; une autre classe encore inhibe les inhibiteurs et ouvre le portillon (proposition~\ref{prop:punctuation}) & Revue des propriétés des classes de neurones \nt{GABA} du cortex. Optogénétique: l'excitation rythmique des neurones \nt{GABA} rapides provoque un rythme gamma, et la réponse à un stimulus dépend de la phase du cycle. Dans le cortex visuel de la souris, les neurones PV s'inhibent mutuellement, les neurones SST inhibent toutes les autres classes, les neurones VIP surtout les SST. L'excitation des neurones VIP chez la souris éveillée lève l'inhibition des neurones \nt{GLUT}; ils sont activés par un signal de renforcement. Le partage des rôles comme règle générale est une hypothèse & \cite{Tremblay2016,Cardin2009,Pfeffer2013,Pi2013} & hypothèse \\
14 & L'énergie limite la fréquence moyenne à environ une impulsion par seconde~\eqref{eq:energy} & Budget énergétique de la substance grise du rongeur d'après des données anatomiques et physiologiques: impulsions et courants synaptiques représentent $47\,\%$ et $34\,\%$ de l'énergie de signalisation; au plus $\sim15\,\%$ des neurones peuvent être actifs simultanément. Pour le cortex humain: au plus $\sim1\,\%$ des neurones peuvent être fortement actifs simultanément. Le nombre de $\sim10^9$ ATP par impulsion et la puissance de $\sim10$\,W pour la signalisation sont des estimations du chapitre fondées sur ces calculs & \cite{Attwell2001,Lennie2003} & théorie \\
15 & Le code est parcimonieux et ajusté pour maximiser le nombre de bits par joule; le cortex échange de la précision contre de l'énergie (proposition~\ref{prop:economy}) & Hypothèse du code économe. Appui: chez des souris soumises à une restriction alimentaire, la conductance des récepteurs AMPA et la dépense synaptique d'ATP baissent dans le cortex visuel ($-29\,\%$), la fréquence des impulsions est conservée, et l'accord sur l'orientation s'élargit de $32\,\%$ & \cite{Barlow1961,Padamsey2022} & hypothèse \\
\bottomrule
\end{longtable}}

\section{Chapitre~\ref{sec:serocomputer}. Les neurones \nt{SERO}}

Les propriétés des neurones \nt{SERO} eux-mêmes (rythme, autoinhibition, signe fixé par le récepteur, sous-systèmes) sont mesurées directement; les calculs du chapitre (régulateur, filtre, logique) découlent de ses équations; $\tau_c$, le coefficient de diffusion de la sérotonine et le nombre de sites de libération sont des estimations; le rôle de la boucle comme régulateur de niveau dans le cerveau est une hypothèse (dans le noyau du raphé, l'autoinhibition est ponctuelle et ne produit que de courtes pauses), et le rôle de \nt{SERO} comme horizon est une hypothèse appuyée par des expériences comportementales.

{\footnotesize
\setlength{\tabcolsep}{3pt}\renewcommand{\arraystretch}{1.15}
\begin{longtable}{>{\raggedright}p{0.5cm}>{\raggedright}p{4.0cm}>{\raggedright}p{5.6cm}>{\raggedright}p{2.3cm}>{\raggedright\arraybackslash}p{1.7cm}}
\toprule
N\textsuperscript{o} & Affirmation & Expérience et ce qui est mesuré & Source & Statut \\
\midrule\endfirsthead
\toprule
N\textsuperscript{o} & Affirmation & Expérience et ce qui est mesuré & Source & Statut \\
\midrule\endhead
1 & Le signe de l'action de la sérotonine est choisi par le récepteur du destinataire (proposition~\ref{prop:receptor}) & Les récepteurs de la sérotonine se répartissent en familles selon leur pharmacologie et leur mécanisme: 5-HT$_{1A}$ ouvre des canaux potassiques par une protéine G et inhibe la cellule, 5-HT$_{2A}$ et le récepteur ionotrope 5-HT$_3$ l'excitent. Un même neurotransmetteur produit des réponses opposées dans des cellules différentes. & \cite{BarnesSharp1999} & mesuré \\
2 & À l'éveil, le neurone \nt{SERO} émet régulièrement quelques impulsions par seconde & Enregistrement de neurones du noyau du raphé dorsal chez le chat libre de ses mouvements: décharge rythmique lente de $2.8\pm0.2$ imp./s à l'éveil calme; en sommeil lent, la fréquence baisse de $34$--$68\,\%$, en sommeil paradoxal de $98\,\%$. Chez le rat anesthésié, $1.7\pm0.5$ imp./s, très régulièrement. & \cite{TrulsonJacobs1979, GagnonParent2014, JacobsAzmitia1992} & mesuré \\
3 & Le neurone \nt{SERO} est un pacemaker: pas besoin d'une poussée pour chaque impulsion, mais il lui faut un apport de fond régulier (en tranche, la noradrénaline par les récepteurs $\alpha_1$; dans l'organisme, \nt{NORA}) & En tranche de cerveau, les cellules déchargent lentement et régulièrement, chaque impulsion étant suivie d'une post-hyperpolarisation de $200$--$800\ms$. Les cellules spontanément actives sont moins nombreuses en tranche que dans l'organisme: le rythme régulier exige une entrée tonique de noradrénaline par les récepteurs $\alpha_1$, dont le blocage l'éteint. & \cite{VandermaelenAghajanian1983} & mesuré \\
4 & Presque tous les récepteurs sont métabotropes; la réponse est lente: elle monte et retombe en une seconde environ & Toutes les familles de récepteurs, sauf 5-HT$_3$, sont couplées à des protéines G. En tranche, une libération évoquée de sérotonine produit par 5-HT$_{1A}$ un courant inhibiteur qui monte et retombe en moins d'une seconde. & \cite{BarnesSharp1999, CourtneyFord2016} & mesuré \\
5 & Dans les régions cibles, la sérotonine est libérée dans le volume, et pas seulement dans la fente; dans le noyau du raphé lui-même, l'inhibition des voisins se fait point à point & Voltampérométrie cyclique rapide en tranche: la libération par impulsion est la même dans une région riche en synapses et dans le noyau du raphé, où il n'y a presque pas de synapses; elle ne dépend pas du blocage de la recapture ni des récepteurs; il y a bien moins de molécules par terminaison que de transporteurs et de récepteurs, et la concentration extracellulaire coïncide avec l'affinité du récepteur principal. Dans le noyau du raphé, l'inhibition par 5-HT$_{1A}$ se fait point à point: le transporteur empêche la sérotonine de s'accumuler hors de la fente. & \cite{Bunin1998, CourtneyFord2016} & mesuré \\
6 & La concentration selon~\eqref{eq:seroneuron} décroît par recapture; $\tau_c$ de l'ordre de la seconde est une estimation & Voltampérométrie dans le cerveau vivant: la sérotonine extracellulaire est limitée avant tout par la recapture et la dégradation, non par la synthèse. Les travaux cités ne mesurent pas $\tau_c$ directement; l'ordre de grandeur est une estimation du chapitre. & \cite{Hashemi2012serotonin} & mesuré; $\tau_c$: estimation \\
7 & NON et NON-OU sur un seul pacemaker; complétude de la logique \nt{SERO} (proposition~\ref{prop:serocomplete}, fig.~\ref{fig:serogates}) & Découle de~\eqref{eq:seroneuron}: un pacemaker inhibable est un NON-OU, et le NON-OU est fonctionnellement complet. Aucune expérience n'a construit de logique avec des neurones \nt{SERO}; la condition de volumes séparés n'est pas remplie dans le cerveau. & démonstration dans le chapitre & théorie \\
8 & La sérotonine inhibe le neurone \nt{SERO} lui-même par des récepteurs qu'il porte & Chez le rat anesthésié, des agonistes 5-HT$_{1A}$ administrés par voie intraveineuse et par iontophorèse inhibent la décharge des neurones du raphé dorsal avec la même relation dose-réponse que la sérotonine elle-même; les agonistes 5-HT$_{1B}$ n'agissent presque pas. En tranche, la sérotonine endogène des cellules voisines produit par 5-HT$_{1A}$ un courant et une courte pause dans la décharge, sans changer la fréquence moyenne. & \cite{Sprouse1987, CourtneyFord2016} & mesuré \\
9 & Dans le modèle, la boucle par un volume commun est un régulateur de niveau: la dispersion et la constante de temps sont divisées par $1+G$~\eqref{eq:feedback}; que la boucle fonctionne ainsi dans le cerveau est une hypothèse & Formule classique de l'amplificateur à rétroaction négative, redémontrée dans le chapitre. Le gain de boucle $G$ du système \nt{SERO} n'est pas mesuré. Ce qui est mesuré: en tranche, l'autoinhibition se fait point à point, produit une courte pause et ne change pas la fréquence moyenne. & \cite{Black1934feedback, CourtneyFord2016} & théorie; dans le cerveau: hypothèse \\
10 & La concentration est une moyenne glissante de la fréquence, un filtre passe-bas~\eqref{eq:lowpass}, fig.~\ref{fig:lowpass} & Solution de l'équation linéaire~\eqref{eq:seroneuron}; la figure est calculée par cette formule pour $\tau_c=1\,\text{s}$. Pas de modèle séparé dans \texttt{models/}. & démonstration dans le chapitre & théorie \\
11 & La tortuosité des espaces extracellulaires ralentit la diffusion d'une petite molécule d'un facteur $2.6$ environ & Méthode de la source ponctuelle: iontophorèse de tétraméthylammonium et enregistrement de sa concentration par une électrode sélective en tranche et dans le cerveau vivant. Tortuosité $\lambda\approx1.6$, c'est-à-dire un $D$ effectif $\lambda^2\approx2.6$ fois plus petit que le $D$ libre; fraction de volume extracellulaire voisine de $0.2$. Le coefficient de diffusion de la sérotonine elle-même dans le tissu n'est pas mesuré dans ces travaux; dans le chapitre, c'est une estimation. & \cite{Sykova2008} & mesuré \\
12 & La longueur d'un \og fil\fg{} indépendant est $\ell\sim\sqrt{D\tau}\approx20$~\textmu m~\eqref{eq:diffusionlength}; le nombre de fils est limité par le nombre de telles régions (proposition~\ref{prop:volumewires}) & Loi de la diffusion et substitution des estimations de $D$ et $\tau$ (lignes~6 et~11); la proposition~\ref{prop:volumewires} en est une conséquence. Aucune expérience directe ne compte les canaux indépendants de la transmission volumique. & démonstration dans le chapitre & théorie \\
13 & La sortie d'un seul neurone \nt{SERO} couvre d'immenses portions du cerveau; $\sim10^5$ sites de libération selon l'estimation & Reconstruction complète de neurones isolés du raphé dorsal chez le rat: tous les axones ascendants se ramifient abondamment, la longueur totale de l'axone d'une cellule atteint $18.7$~cm, les branches d'une même cellule atteignent le striatum, le cortex préfrontal et l'amygdale. Le nombre de sites de libération par cellule n'a pas été compté directement. & \cite{GagnonParent2014} & mesuré; $10^5$: estimation \\
14 & Les neurones \nt{SERO} se répartissent en quelques sous-systèmes aux sorties et aux réponses différentes & Marquage viro-génétique chez la souris: les neurones qui projettent vers l'amygdale et vers le cortex frontal ne se recouvrent presque pas par leurs ramifications, reçoivent des entrées différentes et répondent de façon opposée à un stimulus désagréable; activer ou inactiver chaque groupe modifie le comportement différemment. & \cite{Ren2018} & mesuré \\
15 & Condition de patience $D<\tau_\gamma\ln(R/r_0)$~\eqref{eq:patience} pour une dépréciation exponentielle; pour la dépréciation mesurée, hyperbolique, $D<\tau_\gamma(R/r_0-1)$ & Découle de la dépréciation $e^{-D/\tau_\gamma}$ ou $1/(1+D/\tau_\gamma)$; démonstration dans le chapitre. Singes, choix avec des délais de quelques secondes: la dépréciation est plus proche de l'hyperbole que de l'exponentielle; la réponse des neurones \nt{DOPA} au signal annonçant une récompense différée décroît elle aussi avec le délai selon une hyperbole. & démonstration dans le chapitre; \cite{KobayashiSchultz2008} & théorie \\
16 & \nt{SERO} fixe l'horizon $\tau_\gamma$ (section~\ref{sec:horizon}) & L'activation optogénétique des neurones \nt{SERO} du raphé dorsal chez la souris allonge l'attente d'une récompense différée et accroît la persévérance dans la recherche d'une récompense devenue rare. Chez l'homme, l'abaissement de la sérotonine (régime sans tryptophane) rend la dépréciation d'une récompense différée plus forte. Comment la concentration se transforme en $\tau_\gamma$, on ne le sait pas. & \cite{Doya2002, Miyazaki2014, Lottem2018, Schweighofer2008} & hypothèse \\
\bottomrule
\end{longtable}}

\section{Chapitre~\ref{sec:dofa}. Apprendre avec \nt{DOPA}}

Les mécanismes de la synapse (quanta, détecteur de coïncidences, calcium, fenêtres de plasticité) et le comportement de la dopamine comme erreur de prédiction sont mesurés directement; que les règles mènent au gradient et ce que coûte la diffusion sont des théorèmes; le résultat de l'apprentissage du choix est un calcul par le modèle du livre; le schéma qui calcule l'erreur est une hypothèse.

{\footnotesize
\setlength{\tabcolsep}{3pt}\renewcommand{\arraystretch}{1.15}
\begin{longtable}{>{\raggedright}p{0.5cm}>{\raggedright}p{4.0cm}>{\raggedright}p{5.6cm}>{\raggedright}p{2.3cm}>{\raggedright\arraybackslash}p{1.7cm}}
\toprule
N\textsuperscript{o} & Affirmation & Expérience et ce qui est mesuré & Source & Statut \\
\midrule\endfirsthead
\toprule
N\textsuperscript{o} & Affirmation & Expérience et ce qui est mesuré & Source & Statut \\
\midrule\endhead
1 & La rétropropagation de l'erreur, sous la forme qu'elle a dans les réseaux artificiels, n'a pas été trouvée dans le cerveau & Pas d'expérience directe: c'est une affirmation d'absence. Des revues analysent ce qu'exige l'algorithme (connexions de retour reproduisant les connexions aller, phase séparée) et les approximations biologiques proposées. & \cite{Lillicrap2020, WhittingtonBogacz2019} & hypothèse \\
2 & Le poids se décompose en nombre de sites de libération, probabilité de libération et réponse à un quantum: $w=N_s\,p\,A_1$~\eqref{eq:quantal} & Jonction neuromusculaire de grenouille à libération réduite: la réponse du destinataire varie par paliers multiples de la réponse miniature spontanée à un quantum; le nombre de quanta par impulsion suit une loi de Poisson. & \cite{delCastillo1954} & mesuré \\
3 & Le récepteur du destinataire est un détecteur de coïncidences; le calcium déclenche la modification du poids & Patch-clamp de neurones de souris: les ions Mg$^{2+}$ bouchent au potentiel de repos le canal ouvert par le glutamate et le libèrent lors de la dépolarisation. Tranches d'hippocampe: un chélateur du calcium dans le destinataire bloque la potentialisation à long terme, et la libération photo-induite de calcium dans la cellule renforce à elle seule la transmission. & \cite{Nowak1984, Malenka1988calcium} & mesuré \\
4 & N'importe quel côté exécute la décision: $A_1$ croît chez le destinataire, $p$ change chez l'émetteur & Le glutamate libéré par la lumière au-dessus d'une seule épine provoque une croissance rapide et durable de l'épine et du courant de ses récepteurs; la potentialisation s'accompagne de l'insertion de récepteurs AMPA. La dépolarisation du destinataire supprime temporairement la libération chez l'émetteur; l'effet est porté par des endocannabinoïdes qui retraversent la fente (montré sur des entrées \nt{GABA}). La dépression à court terme dépend de la probabilité de libération de l'émetteur. & \cite{Matsuzaki2004, Malinow2002, WilsonNicoll2001, Tsodyks1997} & mesuré \\
5 & La synapse se renforce quand l'émetteur et le destinataire sont actifs ensemble~\eqref{eq:hebb} & Lapin anesthésié: après de courtes séries d'impulsions sur la voie d'entrée de l'hippocampe, la réponse est renforcée pendant $30$~min à $10$~h. En tranche d'hippocampe, les impulsions du destinataire ne renforcent que les synapses actives au même moment. & \cite{Bliss1973LTP, Kelso1986hebb} & mesuré \\
6 & La règle~\eqref{eq:oja} trouve le vecteur propre principal des corrélations des entrées (proposition~\ref{prop:oja}) & Théorème; démonstration dans le chapitre. Aucune expérience ne montre un neurone ayant appris la composante principale de ses entrées; le mécanisme qui limite la croissance des poids dans la synapse n'est pas établi. & \cite{Oja1982} & théorie \\
7 & Hebb temporel: fenêtre d'environ $\pm20\ms$ (fig.~\ref{fig:stdp}); des séquences répétées font pousser des chaînes & Paires de neurones d'hippocampe en culture: une impulsion du destinataire dans les $20\ms$ qui suivent celle de l'émetteur donne un renforcement durable, dans les $20\ms$ qui précèdent, un affaiblissement; les deux dépendent des récepteurs NMDA. Le rapport $A_-=0.6A_+$ de la figure est une illustration. Que des chaînes poussent ainsi dans un réseau n'a pas été mesuré directement. & \cite{BiPoo1998} & mesuré \\
8 & Trace~\eqref{eq:threefactor}: la dopamine arrivée $1$--$2\,\text{s}$ après l'activité conjointe renforce la connexion, plus tard non (proposition~\ref{prop:threefactor}) & Tranches de striatum, stimulation optogénétique séparée des entrées glutamatergiques et dopaminergiques: l'épine ne grossit que si la dopamine arrive $0.3$--$2\,\text{s}$ après l'entrée glutamatergique; la fenêtre est fixée par la montée rapide et la dégradation de l'AMPc. Dans le cortex, l'appariement laisse des traces séparées pour le renforcement et l'affaiblissement, que les récepteurs des monoamines transforment en modifications durables dans une fenêtre de l'ordre de la seconde. & \cite{Yagishita2014, He2015eligibility} & mesuré \\
9 & Des fenêtres de plusieurs secondes existent aussi sans dopamine: le troisième signal est une forte excitation du destinataire & Souris vivante, champ CA1: un seul événement de plateau dans le destinataire renforce les entrées arrivées quelques secondes avant et après lui, et crée un champ de lieu en un seul passage. & \cite{Bittner2017} & mesuré \\
10 & Le pas moyen de la règle à trois facteurs suit le gradient de la récompense~\eqref{eq:perturbation} (proposition~\ref{prop:gradient}) & Théorème du gradient pour l'apprentissage par écarts aléatoires; dans le chapitre, établi pour un bruit faible. & \cite{Williams1992} & théorie \\
11 & Le bruit est une recherche: le réseau apprend de ses propres écarts aléatoires & Jeunes diamants mandarins: l'inactivation du noyau de sortie du circuit des ganglions de la base supprime brusquement et réversiblement la variabilité du chant. Diamants adultes: une perturbation est appliquée aux exécutions d'une syllabe dont la hauteur est d'un côté de la moyenne, et les oiseaux déplacent vite la hauteur de l'autre côté: ils apprennent d'après leur propre dispersion. & \cite{Olveczky2005, TumerBrainard2007} & mesuré \\
12 & Le choix entre deux options s'apprend avec $\tau_e=1\,\text{s}$ et $\delta=R-\bar R$; pas avec $\tau_e=50\ms$; sans soustraction, les poids croissent sans limite, et avec une grande vitesse d'apprentissage le réseau peut figer le plus mauvais choix & \texttt{models/dofa\_learning.py}, $\eta=0.05$, $500$ essais, $6$ graines. $\tau_e=1\,\text{s}$: part des choix de $A$ $0.65$--$0.79\to0.96$--$1.00$, poids $0.85$--$1.33$. $\tau_e=50\ms$: $0.38$--$0.55$, poids inchangés. $\delta=R$: poids du vainqueur $860$--$1190$. $\delta=R$, $\eta=0.5$, $3$ graines: l'une s'est figée sur $B$ ($0.04\to0$). Le script imprime les quatre cas; le recalcul concorde avec le chapitre. & \texttt{models/\allowbreak dofa\_\allowbreak learning.py} & modèle \\
13 & Le temps d'apprentissage par un seul signal commun croît proportionnellement au nombre de neurones bruités (proposition~\ref{prop:broadcastcost}) & Courbes d'apprentissage d'un réseau linéaire: la perturbation des neurones est plus lente que le gradient exact d'un facteur égal au nombre de neurones de sortie, la perturbation des poids encore d'un facteur égal au nombre d'entrées. & \cite{Werfel2005} & théorie \\
14 & Les sorties de \nt{DOPA} vont avant tout aux régions qui choisissent les actions & Reconstruction complète des axones de neurones \nt{DOPA} nigrostriés isolés chez le rat: les branches d'une cellule couvrent $0.45$--$5.7\,\%$ du volume du striatum (en moyenne $2.7\,\%$), et il n'y a presque pas de branches hors du striatum. & \cite{Matsuda2009} & mesuré \\
15 & Le cortex apprend à prédire son entrée, l'erreur est calculée sur place & Revue: le cortex sensoriel présente des réponses au désaccord entre l'entrée attendue et l'entrée reçue. Qu'il s'agisse d'une règle générale d'apprentissage du cortex n'est pas démontré. & \cite{Keller2018} & hypothèse \\
16 & La dopamine se comporte comme l'erreur~\eqref{eq:ctd}: pic, transfert sur le signal annonciateur, creux (fig.~\ref{fig:ctdcases}) & Neurones \nt{DOPA} du singe: pic pour un jus inattendu; après apprentissage, pic sur le signal annonciateur et pas de réponse au jus promis; creux quand le jus n'arrive pas. La forme continue~\eqref{eq:ctd} relève de la théorie. & \cite{Schultz1997, Doya2000} & mesuré \\
17 & Schéma qui calcule $\dd V/\dd t$ et soustrait l'attente & Souris, enregistrement et optogénétique dans l'aire tegmentale ventrale: les neurones \nt{DOPA} soustraient l'attente de la récompense; les neurones \nt{GABA} voisins les inhibent quand la récompense est attendue, et les exciter ou les inhiber modifie l'erreur. Le calcul de la dérivée n'a pas été montré. & \cite{Eshel2015arithmetic} & hypothèse \\
18 & Le neurone \nt{DOPA} est un pacemaker; les creux sont moins profonds que les pics & En tranche, les neurones \nt{DOPA} déchargent spontanément et très régulièrement, sur un potentiel pacemaker lent. Chez le singe, la fréquence suit quantitativement l'erreur positive, mais ne transmet que faiblement l'erreur négative. & \cite{GraceOnn1989, Bayer2005} & mesuré \\
19 & L'acteur et le critique (fig.~\ref{fig:actorcritic}) & L'algorithme vient de la théorie de l'apprentissage par renforcement. Chez l'homme (IRMf), le striatum ventral se comporte comme un critique avec ou sans choix, le striatum dorsal seulement quand il faut choisir des actions. & \cite{SuttonBarto2018, ODoherty2004actor} & hypothèse \\
20 & Le signe de $\delta$ dans la règle est inversé chez les neurones de la seconde famille de récepteurs & Tranches de striatum: chez les neurones à récepteurs D1, l'appariement produit un renforcement qui exige D1; chez les neurones à D2, un affaiblissement qui exige D2. & \cite{Shen2008} & mesuré \\
\bottomrule
\end{longtable}}

\section{Chapitre~\ref{sec:dofaalt}. Autres théories de \nt{DOPA}}

Chacune des images élargies s'appuie sur ses propres mesures directes de la dopamine (dispersion des points d'inversion, réponses au désagréable et au nouveau, montée lente, réponses au changement de goût), mais les formules qui les expliquent sont des théories, et pour deux d'entre elles les expériences se contredisent encore.

{\footnotesize
\setlength{\tabcolsep}{3pt}\renewcommand{\arraystretch}{1.15}
\begin{longtable}{>{\raggedright}p{0.5cm}>{\raggedright}p{4.0cm}>{\raggedright}p{5.6cm}>{\raggedright}p{2.3cm}>{\raggedright\arraybackslash}p{1.7cm}}
\toprule
N\textsuperscript{o} & Affirmation & Expérience et ce qui est mesuré & Source & Statut \\
\midrule\endfirsthead
\toprule
N\textsuperscript{o} & Affirmation & Expérience et ce qui est mesuré & Source & Statut \\
\midrule\endhead
1 & La règle asymétrique~\eqref{eq:asymmetric} converge vers le point~\eqref{eq:expectile}; pour $\kappa=1/2$ c'est la moyenne, pour une goutte de probabilité $1/2$, $V=\kappa$ (fig.~\ref{fig:expectiles}) & Le point~\eqref{eq:expectile} est un expectile, solution d'un problème de moindres carrés avec des poids différents pour les écarts positifs et négatifs; pour l'exemple du chapitre, démonstration dans le chapitre même. & \cite{NeweyPowell1987}; démonstration dans le chapitre & théorie \\
2 & Les différents neurones \nt{DOPA} ont des points d'inversion différents, ordonnés selon l'asymétrie (proposition~\ref{prop:expectile}) & Souris, neurones \nt{DOPA} de l'aire tegmentale ventrale marqués par optogénétique, récompenses de tailles différentes: le point où le pic fait place au creux diffère selon les neurones; il est plus haut chez les neurones aux réponses plus asymétriques; les réponses du groupe permettent de reconstituer la forme de la distribution des récompenses. Que ce soit la fonction principale de la dispersion est une hypothèse. & \cite{Dabney2020} & mesuré \\
3 & Une partie des neurones \nt{DOPA} répond au désagréable par un pic & Singes: certains neurones \nt{DOPA} sont excités par l'indice d'une récompense et inhibés par l'indice d'un jet d'air dans l'œil, d'autres sont excités par les deux; les groupes se trouvent dans des parties différentes du mésencéphale. & \cite{Matsumoto2009, BrombergMartin2010} & mesuré \\
4 & Le module de l'erreur ou la nouveauté fixe la vitesse d'apprentissage & Les travaux cités ne contiennent pas d'expérience directe où le signal d'importance des neurones \nt{DOPA} change la vitesse d'apprentissage; la revue propose un rôle de signal \og attention, c'est important\fg{}. & \cite{BrombergMartin2010} & hypothèse \\
5 & Un fil séparé pour l'axe du danger & Souris: les neurones \nt{DOPA} qui projettent vers l'extrémité postérieure du striatum répondent aux stimuli nouveaux et intenses, et non à la récompense; leur destruction affaiblit l'évitement d'un objet nouveau. & \cite{Menegas2018} & mesuré \\
6 & Temps d'action optimal $\tau_a^*=\sqrt{C/\bar R}$~\eqref{eq:vigor} & Minimum de la somme $C/\tau_a+\bar R\tau_a$; démonstration dans le chapitre. Dans le modèle d'origine, la vitesse des actions est fixée par le prix du temps, égal à la récompense moyenne. & \cite{Niv2007}; démonstration dans le chapitre & théorie \\
7 & Le niveau lent de dopamine porte $\bar R$ et fixe la vitesse des actions (proposition~\ref{prop:multiplex}) & Rats, microdialyse et voltampérométrie dans le noyau accumbens: le niveau de dopamine suit, de minute en minute, la fréquence des récompenses et la vitesse des actions. Chez l'homme, la L-DOPA renforce la dépendance du temps de réaction à la fréquence moyenne des récompenses. Que le niveau lent soit précisément $\bar R$ n'est pas démontré. & \cite{Hamid2016, Beierholm2013vigor} & hypothèse \\
8 & Le niveau lent a deux sources: la libération varie aux sorties sans que change la fréquence des impulsions, mais une montée lente existe aussi dans la fréquence elle-même & Rats, même tâche: la libération dans le noyau accumbens suit l'attente changeante de la récompense, mais pas la fréquence des impulsions des neurones \nt{DOPA} identifiés de l'aire tegmentale ventrale. Souris dans un couloir virtuel (ligne~10): la montée progressive à l'approche de la récompense existe aussi dans les impulsions de neurones \nt{DOPA} identifiés. & \cite{Mohebi2019, Kim2020} & mesuré \\
9 & La dopamine monte progressivement pendant que l'animal s'approche d'une récompense lointaine & Rats dans un labyrinthe, voltampérométrie dans le striatum: la concentration de dopamine augmente en quelques secondes à mesure que le but approche; la pente dépend de la distance et de la taille de la récompense. & \cite{Howe2013ramp, Hamid2016} & mesuré \\
10 & La montée est l'erreur $\delta$ quand l'estimation se précise~\eqref{eq:ramp}; un saut dans le couloir donne un pic (fig.~\ref{fig:ramp}) & La formule et la valeur $\delta=0.75\,V$ par seconde sont démontrées dans le chapitre. Expérience: souris dans un couloir virtuel, téléportation instantanée plus près du but; les impulsions des corps cellulaires, le calcium des corps et des terminaisons et la concentration dans le striatum répondent comme une erreur, et non comme l'attente $V$. Laquelle des deux explications vaut pour toutes les sorties est en discussion. & \cite{Kim2020} & mesuré \\
11 & Regard en arrière: la dopamine signale la mesure du lien causal~\eqref{eq:retro} & Souris, libération de dopamine dans le noyau accumbens dans des expériences où cette théorie et l'erreur~\eqref{eq:ctd} prédisent des choses différentes: dans plusieurs expériences, la libération a suivi la première. & \cite{Jeong2022} & hypothèse \\
12 & Au cours de l'apprentissage, la réponse dopaminergique se déplace progressivement de la récompense vers l'indice & Souris, signal des terminaisons des neurones \nt{DOPA} dans le striatum ventral au fil de l'apprentissage: le pic passe progressivement du moment de la récompense au moment de l'indice, comme l'exige l'apprentissage selon~\eqref{eq:ctd}. & \cite{Amo2022} & mesuré \\
13 & Les neurones \nt{DOPA} répondent au changement de goût de la récompense à valeur égale & Rats, enregistrement de neurones de l'aire tegmentale ventrale: le remplacement d'un jus par un autre de même valeur provoque une réponse, bien que l'erreur~\eqref{eq:ctd} soit nulle. & \cite{Takahashi2017} & mesuré \\
14 & Des pics artificiels enseignent le lien entre deux stimuli neutres & Rats, expérience de préconditionnement sensoriel sans récompense: l'activation optogénétique des neurones \nt{DOPA} lors du passage du premier stimulus au second crée le lien, leur inactivation empêche de l'apprendre. & \cite{Sharpe2017} & mesuré \\
15 & Vecteur d'erreurs par caractéristique~\eqref{eq:vectortd} & Théorie de l'erreur de prédiction des caractéristiques; la forme vectorielle n'a pas été testée directement. & \cite{Gardner2018} & hypothèse \\
16 & Dans une même tâche, différents neurones \nt{DOPA} portent différentes grandeurs & Souris dans un couloir virtuel, imagerie à deux photons des neurones \nt{DOPA}: certains sont liés à la récompense, d'autres à la vitesse et aux virages, d'autres encore aux indices et à la justesse du choix. & \cite{Engelhard2019} & mesuré \\
17 & Un faisceau au lieu d'un fil (proposition~\ref{prop:bundle}) & Synthèse des lignes~2--16: chaque décomposition est mesurée séparément; aucune expérience ne donne une image unifiée où elles seraient toutes les parties d'un même signal. & --- & hypothèse \\
\bottomrule
\end{longtable}}

\section{Chapitre~\ref{sec:nora}. \nt{NORA}}

L'organisation du système \nt{NORA}, ses deux modes, le lien avec la pupille (qui ne passe pas seulement par \nt{NORA}), l'augmentation du rapport du signal au fond et le U inversé du comportement sont mesurés directement; le gain multiplicatif $g$ est un modèle; que le gain change le régime du choix et agisse comme l'inverse d'une température sont des conclusions du chapitre; le bénéfice de l'ajustement du gain est un calcul par le modèle du livre; \og \nt{NORA}, bouton de l'exploration\fg{} et \og \nt{NORA}, interruption\fg{} sont des hypothèses.

{\footnotesize
\setlength{\tabcolsep}{3pt}\renewcommand{\arraystretch}{1.15}
\begin{longtable}{>{\raggedright}p{0.5cm}>{\raggedright}p{4.0cm}>{\raggedright}p{5.6cm}>{\raggedright}p{2.3cm}>{\raggedright\arraybackslash}p{1.7cm}}
\toprule
N\textsuperscript{o} & Affirmation & Expérience et ce qui est mesuré & Source & Statut \\
\midrule\endfirsthead
\toprule
N\textsuperscript{o} & Affirmation & Expérience et ce qui est mesuré & Source & Statut \\
\midrule\endhead
1 & L'homme possède quelques dizaines de milliers de neurones \nt{NORA}, presque tous en un seul groupe & Coupes sériées du tronc cérébral humain, marquage de la tyrosine hydroxylase: $53\,900$ cellules dans le locus coeruleus et encore $6\,260$ dans un sous-noyau voisin. & \cite{Baker1989LC} & mesuré \\
2 & Les sorties se répartissent dans presque tout le cerveau, mais les parties du groupe desservent en partie des régions différentes & Marquage viro-génétique chez la souris: les neurones \nt{NORA} du locus coeruleus reçoivent des entrées de vastes régions du cerveau et envoient des sorties vers de vastes régions du cerveau et de la moelle épinière. Chez le rat: le groupe qui projette vers l'amygdale renforce l'apprentissage de la peur, et un groupe distinct projetant vers le cortex préfrontal renforce son extinction. & \cite{Schwarz2015LC, Uematsu2017LC, Poe2020} & mesuré \\
3 & Mode tonique: impulsions régulières, d'une fraction d'impulsion à quelques impulsions par seconde, fréquence qui varie avec l'éveil & Rats, enregistrement de neurones du locus coeruleus pendant le cycle veille-sommeil: la fréquence est maximale à l'éveil, plus basse en sommeil lent, presque nulle en sommeil paradoxal; les variations de fréquence précèdent le changement d'état. & \cite{AstonJonesBloom1981} & mesuré \\
4 & Mode phasique: courte bouffée sur un événement important pour la tâche, à peu près au moment de la décision & Singes, tâche de détection d'une cible rare: tous les neurones du locus coeruleus répondent par une bouffée à la cible seulement, $\sim90\ms$ après elle et $\sim200\ms$ avant la réponse manuelle; quand la performance est mauvaise, les bouffées sont plus faibles. Dans une tâche de choix entre deux options, la bouffée est plus étroitement liée à la réponse qu'au stimulus, et elle est présente pour les réponses justes comme pour les erreurs. & \cite{AstonJones1994vigilance, Clayton2004LC} & mesuré \\
5 & Le diamètre de la pupille est un point de mesure imprécis de \nt{NORA}: il suit \nt{NORA}, mais aussi \nt{ACET} et d'autres régions & Singes: les impulsions du locus coeruleus, jusqu'aux impulsions isolées d'une seule cellule, prédisent la dilatation de la pupille; un lien semblable, mais plus tardif, existe dans les colliculus et le cortex cingulaire. Souris: les fluctuations de la pupille suivent l'activité des sorties noradrénergiques comme cholinergiques dans le cortex. & \cite{Joshi2016, Reimer2016} & mesuré \\
6 & La noradrénaline supprime l'activité de fond plus que l'activité évoquée; le gain multiplicatif~\eqref{eq:gain} est un modèle qui rend compte de cet effet & Singes éveillés, cortex auditif, iontophorèse de noradrénaline: l'activité de fond des neurones est plus supprimée que la réponse aux vocalisations des congénères; le rapport signal sur bruit augmente. La sigmoïde multiplicative~\eqref{eq:gain} est reprise d'un modèle de réseau; la multiplication littérale de la pente n'est pas mesurée. & \cite{Foote1975NE, ServanSchreiber1990} & mesuré; $g$: modèle \\
7 & Le gain n'augmente $d'$ que contre le bruit postérieur à l'amplificateur~\eqref{eq:gainsnr} (proposition~\ref{prop:gainnoise}) & Démonstration dans le chapitre; dans un modèle de réseau, le gain améliore la détection du signal. Aucune expérience ne sépare le bruit avant et après l'amplificateur pour tester cette distinction. & démonstration dans le chapitre; \cite{ServanSchreiber1990} & théorie \\
8 & La frontière $g\beta=1$ fait passer le réseau de choix de \og réfléchit\fg{} à \og a décidé\fg{}~\eqref{eq:wtagain} (proposition~\ref{prop:gainwta}, fig.~\ref{fig:pitchfork}) & Analyse linéaire de deux groupes qui s'inhibent mutuellement; démonstration dans le chapitre. Ce seuil n'a pas été mesuré directement dans le cerveau. & démonstration dans le chapitre & théorie \\
9 & Le gain est l'inverse de la température du choix~\eqref{eq:softmax} & Avec un bruit logistique après l'amplificateur, la probabilité de choix est un softmax avec $T=\sigma/g$; démonstration dans le chapitre. & démonstration dans le chapitre & théorie \\
10 & \nt{NORA} fixe combien chercher: forte activité tonique, petit $g$ effectif (exploration); bouffée phasique au moment de la décision, grand $g$ (décision) & Humains: avant un choix au hasard, la pupille de base est plus large qu'avant le choix de la meilleure option connue; les personnes à pupille plus large explorent davantage. Rats: le passage au mode de choix aléatoire est commandé par l'entrée \nt{NORA} dans le cortex cingulaire antérieur. Que la bouffée augmente le $g$ effectif n'est pas mesuré directement. & \cite{JepmaNieuwenhuis2011, Tervo2014stochastic, AstonJones2005} & hypothèse \\
11 & La récompense dépend de $g$ comme un U inversé; le meilleur $g$ est plus petit dans un monde qui change vite (proposition~\ref{prop:explore}, fig.~\ref{fig:invertedU}) & \texttt{models/nora\_explore.py}, $60$ simulations de $4000$ essais. Changement tous les $1000$ essais: meilleur $g=16$, part $0.976$; tous les $20$: $g=8$, part $0.886$; pour $g=0.5$, $0.77$. Le recalcul concorde avec le chapitre. & \texttt{models/\allowbreak nora\_\allowbreak explore.py} & modèle \\
12 & U inversé dans le comportement: meilleur résultat avec une activité tonique modérée et des bouffées nettes & Singes, même tâche qu'à la ligne~4: pendant les périodes de bonne performance, la fréquence tonique est modérée et les bouffées sur la cible sont nettes; avec une fréquence tonique élevée, il n'y a pas de bouffées et les fausses réponses sont plus nombreuses. Les variations de l'activité tonique et évoquée suivent les fluctuations de la qualité de la performance. & \cite{Usher1999LC, AstonJones2005} & mesuré \\
13 & \nt{NORA} signale l'incertitude inattendue, \nt{ACET} l'incertitude attendue & Théorie de l'inférence bayésienne à deux sortes d'incertitude; les travaux cités ne contiennent pas d'expérience directe séparant ces signaux selon le neurotransmetteur. & \cite{YuDayan2005} & hypothèse \\
14 & Après un changement brusque, les humains apprennent plus vite, et la pupille se dilate & Humains, tâche de prédiction avec changements soudains: les variations brèves de la pupille suivent l'estimation de la probabilité d'un changement et, avec la pupille de base, prédisent à quel point les nouvelles données modifient l'estimation (vitesse d'apprentissage); une variation de la pupille sans lien avec la lumière ni avec la tâche modifie cette vitesse. Que la pupille indique ici précisément \nt{NORA} est une conclusion tirée des données indirectes de la ligne~5. & \cite{Nassar2012} & mesuré \\
15 & La bouffée phasique fonctionne comme une interruption: le réseau remet son état à zéro & Aucune expérience directe ne montre une bouffée \nt{NORA} remettant à zéro l'état du réseau; seules les bouffées sur les événements importants sont mesurées (ligne~4). & \cite{BouretSara2005} & hypothèse \\
16 & L'ajustement de $g$ ou de la vitesse d'apprentissage par la surprise ne fait guère mieux que les meilleurs réglages constants & \texttt{models/nora\_explore.py}, monde à époques de $1000$ essais (changement tous les $1000$ et tous les $20$ essais). Meilleur $g$ constant, $g=8$: $0.925$; ajustement de $g$: jusqu'à $0.927$; ajustement de la vitesse d'apprentissage: jusqu'à $0.926$; vitesse constante $0.4$: $0.928$. Le recalcul concorde avec le chapitre. & \texttt{models/\allowbreak nora\_\allowbreak explore.py} & modèle \\
\bottomrule
\end{longtable}}

\section{Chapitre~\ref{sec:acet}. Ajoutons \nt{ACET}}

Les trois actions de l'acétylcholine sont mesurées sur tranches et dans le cerveau vivant, le conflit entre écriture et lecture est montré sur le modèle du livre, et que \nt{ACET} serve de ligne de validation d'écriture et signale l'incertitude attendue est une hypothèse en partie appuyée par l'expérience.

{\footnotesize
\setlength{\tabcolsep}{3pt}\renewcommand{\arraystretch}{1.15}
\begin{longtable}{>{\raggedright}p{0.5cm}>{\raggedright}p{4.0cm}>{\raggedright}p{5.6cm}>{\raggedright}p{2.3cm}>{\raggedright\arraybackslash}p{1.7cm}}
\toprule
N\textsuperscript{o} & Affirmation & Expérience et ce qui est mesuré & Source & Statut \\
\midrule\endfirsthead
\toprule
N\textsuperscript{o} & Affirmation & Expérience et ce qui est mesuré & Source & Statut \\
\midrule\endhead
1 & Les neurones \nt{ACET} du cerveau sont peu nombreux, réunis en quelques groupes, et leurs sorties se répartissent dans tout le cortex: c'est un canal de diffusion & Marquage par anticorps de l'enzyme de synthèse de l'acétylcholine et traçage rétrograde chez le rat: le cortex, l'hippocampe, l'amygdale, le bulbe olfactif et le thalamus reçoivent l'acétylcholine surtout de six groupes compacts de cellules du télencéphale basal et du tronc cérébral & \cite{Mesulam1983} & mesuré \\
2 & Le niveau d'acétylcholine dans le cortex est élevé à l'éveil et bas en sommeil lent & Microdialyse dans le cortex et l'hippocampe de chats libres de leurs mouvements, avec enregistrement EEG: la libération d'acétylcholine est supérieure de $\sim100\%$ à l'éveil calme et de $\sim175\%$ à l'éveil actif par rapport au sommeil lent; en sommeil paradoxal, dans le cortex, elle est la même qu'à l'éveil & \cite{Marrosu1995,Hasselmo1999} & mesuré \\
3 & Le sens du signal est fixé par le récepteur: l'acétylcholine a des récepteurs-canaux rapides et des récepteurs lents qui déclenchent des réactions dans la cellule (proposition~\ref{prop:receptor}) & Revue de mesures: l'action de l'acétylcholine sur l'excitabilité, la transmission et la plasticité dépend du lieu de libération, du sous-type de récepteur (nicotiniques: canaux ioniques; muscariniques: par des protéines G) et de la cellule destinataire & \cite{Picciotto2012} & mesuré \\
4 & Première action: affaiblir les connexions internes sans presque toucher aux entrées extérieures (facteur $1-a$ dans~\eqref{eq:acetnet}) & Tranches de cortex olfactif de rat: à $100$\,\textmu M d'agonistes de l'acétylcholine, les potentiels des connexions internes (couche Ib) baissent de plus de $60\%$, ceux de l'entrée externe (couche Ia) de moins de $15\%$, dans la même cellule; la suppression est présynaptique. Tranches d'hippocampe (CA1): $89\%$ contre $40\%$ pour la voie interne et la voie d'entrée & \cite{HasselmoBower1992,HasselmoSchnell1994} & mesuré \\
5 & Deuxième action: renforcer l'entrée (facteur $\frac{1+a}{2}$) & Tranches de néocortex: les récepteurs nicotiniques renforcent les synapses de l'entrée thalamique et n'agissent pas sur les synapses intracorticales; les muscariniques affaiblissent les deux types. L'acétylcholine augmente l'excitabilité des neurones (revue) & \cite{Gil1997,Hasselmo2006} & mesuré \\
6 & Troisième action: faciliter la modification des poids (vitesse $\eta a$) & Rat: la stimulation électrique du noyau basal (source de l'acétylcholine du cortex), appariée à un son, provoque chez l'animal adulte une forte réorganisation de la carte du cortex auditif en faveur du son présenté & \cite{KilgardMerzenich1998,Hasselmo2006} & mesuré \\
7 & La règle de Hebb pour motifs clairsemés, dans la seconde équation~\eqref{eq:acetnet}, donne une mémoire associative qui complète un motif à partir d'une partie & Calcul de la capacité d'un réseau à motifs clairsemés et règle de covariance: pour une faible fraction de neurones actifs $p$, le réseau stocke nettement plus de motifs que pour $p=1/2$ & \cite{Tsodyks1988} & théorie \\
8 & Écrire avec $a$ bas abîme la mémoire: le réseau écrit ce dont il se souvient, et non ce qu'il voit; pour $a=0.1$, les dégâts ne sont pas moindres que pour $a=0$ & \texttt{models/\allowbreak acet\_memory.py}: $200$ neurones, cinq motifs de $20$, un nouveau motif partage $10$ neurones avec l'un d'eux, $10$ simulations. Pour $a=0$, le nouveau motif n'est pas mémorisé, l'ancien entraîne $5.9$ neurones étrangers sur $10$; pour $a=0.1$, le nouveau motif est rappelé ($0.97$), mais les deux fusionnent: le nouveau entraîne $7.3$ neurones étrangers, l'ancien $7.9$; à partir de $a=0.2$, moins d'un & démonstration dans le chapitre & modèle \\
9 & Proposition~\ref{prop:writeread}: aucun niveau $a$ ne convient aussi bien à l'écriture qu'à la lecture (fig.~\ref{fig:acetsim}) & Même script, vitesse d'apprentissage $\eta a$: le nouveau motif est entièrement mémorisé en une présentation pour $a\ge0.9$, à $0.8$ pour $a=0.75$, pas du tout pour $a\le0.5$. Lecture d'après la moitié du motif: $1.0$ pour $a\le0.2$, $0.96$ pour $a=0.3$, $0.04$ pour $a=0.4$ & démonstration dans le chapitre & modèle \\
10 & Dans le cerveau, un seul fil de diffusion, \nt{ACET}, commute écriture et lecture & L'idée vient de modèles construits à partir de mesures sur tranches. L'expérience la plus directe est humaine: la scopolamine, qui bloque les récepteurs muscariniques, gêne l'apprentissage de nouvelles paires de mots, surtout celles qui recoupent d'anciennes, mais pas le rappel des paires déjà apprises. Les deux modes du réseau n'ont pas été mesurés directement & \cite{HasselmoAndersonBower1992,Atri2004} & hypothèse \\
11 & Le filtre~\eqref{eq:kalman} est optimal; le poids de l'entrée et la vitesse d'apprentissage sont une même grandeur $P/R$; en régime établi, $P=\sqrt{qR}$ et la mémoire vaut $\sqrt{R/q}$ (proposition~\ref{prop:kalman}) & Aucune expérience n'est requise: l'optimalité du filtre de Kalman--Bucy est un résultat classique de la théorie de l'estimation; le régime établi est démontré dans le chapitre & \cite{KalmanBucy1961} & théorie \\
12 & \nt{ACET} communique au réseau une grandeur semblable à $P$, l'incertitude attendue & Proposé dans un modèle probabiliste de l'attention. Il n'existe pas de mesure directe montrant que le niveau d'acétylcholine suit l'incertitude de l'estimation & \cite{YuDayan2005} & hypothèse \\
13 & Une partie des neurones \nt{ACET} répond à la récompense et à la punition en une vingtaine de millisecondes, d'autant plus fort que le renforcement est inattendu & Souris, neurones cholinergiques du télencéphale basal identifiés par optogénétique: réponse à la récompense et à la punition après $18\pm3$\,ms, amplitude croissant avec le caractère inattendu du renforcement, réponses semblables chez les neurones de deux noyaux & \cite{Hangya2015} & mesuré \\
14 & \nt{NORA} remarque un changement inattendu du monde, \nt{ACET} maintient le réseau en mode écriture & Division du travail proposée dans un modèle. Indirectement: chez l'homme, le diamètre de la pupille, lié à \nt{NORA}, suit les changements et la vitesse d'apprentissage. Le couple \nt{NORA}--\nt{ACET} n'a pas été testé directement & \cite{YuDayan2005,Nassar2012} & hypothèse \\
15 & En sommeil lent, le réseau en mode lecture rejoue ce qui a été écrit pendant la journée, et ces répétitions le recopient dans la mémoire à long terme & Enregistrements d'ensembles de neurones de l'hippocampe du rat: les cellules qui ont travaillé ensemble pendant la journée se déclenchent plus souvent ensemble pendant le sommeil lent qui suit, et dans le même ordre; c'est mesuré. Pour la recopie: supprimer les ondulations hippocampiques (ripples) après l'apprentissage dégrade la mémoire spatiale, et allonger les ondulations spontanées l'améliore; que la cause soit le bas niveau d'\nt{ACET} n'est pas démontré & \cite{WilsonMcNaughton1994,Skaggs1996,Girardeau2009,FernandezRuiz2019,Hasselmo1999} & hypothèse \\
\bottomrule
\end{longtable}}

\section{Chapitre~\ref{sec:synthesis}. Toute la machine}

Ce chapitre de synthèse n'apporte pas d'expériences nouvelles: chaque ligne s'appuie sur le chapitre où l'affirmation est examinée et sur les mêmes sources; le bus \nt{GLUT} et la commande du flux par \nt{GABA} sont solidement établis, tandis que les rôles des canaux de diffusion et leur fonctionnement conjoint sont pour l'essentiel une hypothèse.

{\footnotesize
\setlength{\tabcolsep}{3pt}\renewcommand{\arraystretch}{1.15}
\begin{longtable}{>{\raggedright}p{0.5cm}>{\raggedright}p{4.0cm}>{\raggedright}p{5.6cm}>{\raggedright}p{2.3cm}>{\raggedright\arraybackslash}p{1.7cm}}
\toprule
N\textsuperscript{o} & Affirmation & Expérience et ce qui est mesuré & Source & Statut \\
\midrule\endfirsthead
\toprule
N\textsuperscript{o} & Affirmation & Expérience et ce qui est mesuré & Source & Statut \\
\midrule\endhead
1 & Pour $8.6\cdot10^{10}$ neurones, il n'existe que quatre sortes de signaux de commande~\eqref{eq:machine} & Comptage des noyaux cellulaires dans un homogénat de cerveau (fractionneur isotrope): $86.1\pm8.1$ milliards de neurones chez l'homme adulte & \cite{Azevedo2009} & mesuré \\
2 & Proposition~\ref{prop:architecture}, deux premières couches: les données circulent sur le bus d'adresses \nt{GLUT}, le signe de la connexion est fixé par l'émetteur, les neurones \nt{GABA} orientent et normalisent le flux (chapitres~\ref{sec:neurontypes}, \ref{sec:gaba}) & Un seul neurotransmetteur principal par neurone (revue des mesures); l'inhibition directe rétrécit la fenêtre dans laquelle le neurone pyramidal additionne ses entrées; la division par l'activité totale est mesurée dans de nombreuses régions & \cite{Strata1999,Pouille2001,Carandini2012} & mesuré \\
3 & Les éléments ne sont pas fiables: une synapse perd la plupart des impulsions (tabl.~\ref{tab:computer}, chapitre~\ref{sec:code}) & Tranches d'hippocampe, stimulation minimale: plus de la moitié des impulsions ne provoquent aucune réponse dans CA1; les échecs du seuil et de la conduction axonale sont exclus, la cause est la libération probabiliste du neurotransmetteur & \cite{AllenStevens1994} & mesuré \\
4 & Le cerveau fonctionne avec $\sim20$ watts, en moyenne environ une impulsion par seconde et par neurone (chapitre~\ref{sec:code}); les ingénieurs n'ont pas encore construit une telle machine & Estimation~\eqref{eq:energy} à partir des coûts mesurés: environ $10^9$ molécules d'ATP par impulsion, travail des synapses compris (cortex de rongeur); une estimation détaillée pour le cortex donne une fréquence encore plus basse. État de la technique d'après une revue & \cite{Attwell2001,Lennie2003,MehonicKenyon2022} & théorie \\
5 & L'apprentissage de la plupart des synapses est le produit d'une trace locale et d'un seul nombre commun $\delta$ (deuxième équation de~\eqref{eq:machine}, chapitre~\ref{sec:dofa}) & Les neurones \nt{DOPA} du singe répondent à l'erreur de prédiction de la récompense; dans des tranches de striatum, la dopamine n'agrandit une épine que si elle arrive $0.3$--$2$\,s après l'entrée glutamatergique: la trace est mesurée & \cite{Schultz1997,Yagishita2014} & mesuré \\
6 & Un fil d'erreur propre à chaque sortie n'existe que dans le circuit d'apprentissage moteur (chapitre~\ref{sec:motorlearning}) & Cervelet: la coïncidence du signal de la fibre grimpante avec une entrée affaiblit cette entrée; chez le singe, les impulsions complexes des cellules de Purkinje portent la direction de l'erreur de mouvement et déclenchent une correction & \cite{Ito2001,Herzfeld2018} & mesuré \\
7 & Chaque canal de diffusion est un faisceau de quelques lignes (proposition~\ref{prop:bundle}) & Pour \nt{DOPA}: dans une même tâche, des neurones différents portent la récompense, le mouvement et les propriétés du stimulus; l'erreur rapide et le niveau lent de dopamine sont distincts. Pour \nt{SERO}, \nt{NORA}, \nt{ACET}, cette décomposition est moins étudiée & \cite{Engelhard2019,Hamid2016,Mohebi2019} & mesuré \\
8 & \nt{SERO}: régulateur de niveau et, par hypothèse, horizon $\tau_\gamma$ (chapitre~\ref{sec:serocomputer}) & Auto-inhibition: les agonistes de leurs propres récepteurs 5-HT$_{1A}$ inhibent les neurones sérotoninergiques du raphé (mesuré). L'horizon est proposé en théorie; l'expérience la plus forte: l'activation optogénétique de ces neurones chez la souris allonge l'attente d'une récompense différée & \cite{Sprouse1987,Doya2002,Miyazaki2014} & hypothèse \\
9 & \nt{NORA}: le gain $g$ choisit entre chercher et décider (chapitre~\ref{sec:nora}) & Hausse du rapport signal/bruit: chez le singe éveillé, la noradrénaline dans le cortex auditif supprime l'activité de fond plus fortement que la réponse aux vocalisations des congénères (mesuré); le gain multiplicatif est un modèle; le rôle dans le choix entre chercher et décider est une théorie & \cite{Foote1975NE,ServanSchreiber1990,AstonJones2005} & hypothèse \\
10 & \nt{ACET}: bascule entre écriture et lecture (chapitre~\ref{sec:acet}) & Les trois actions (affaiblissement des connexions internes, renforcement de l'entrée, facilitation de la plasticité) sont mesurées; la bascule des modes est soutenue indirectement par une expérience à la scopolamine chez l'homme & \cite{HasselmoBower1992,Gil1997,Atri2004} & hypothèse \\
11 & La formule~\eqref{eq:machine} décrit la machine entière & C'est un résumé des modèles des chapitres~\ref{sec:glutcomputer}--\ref{sec:acet}, chaque partie s'appuie sur son chapitre; l'ensemble n'a pas été vérifié par l'expérience & déduction du chapitre & hypothèse \\
12 & Les boutons fonctionnent de concert sans commande centrale: erreur, surprise, écriture, retour (fig.~\ref{fig:timeline}) & Pas d'expérience: le schéma est qualitatif. Maillons isolés: théorie de la division du travail entre \nt{NORA} et \nt{ACET}, et lien entre la pupille et la vitesse d'apprentissage chez l'homme & \cite{YuDayan2005,Nassar2012} & hypothèse \\
13 & L'apprentissage par un seul signal commun est d'autant plus lent que plus de neurones influent sur le résultat (proposition~\ref{prop:broadcastcost}) & Calcul des courbes d'apprentissage d'un réseau linéaire: l'apprentissage par perturbation des sorties est plus lent que la descente de gradient d'un facteur égal au nombre de sorties & \cite{Werfel2005} & théorie \\
14 & On ignore comment le cortex règle ses circuits multicouches; candidats: bouffées d'impulsions, calculs dans les branches du neurone, apprentissage auto-supervisé par prédiction & Bouffées comme porteuses de l'erreur: modèle; seul un réseau de $5$--$8$ couches a su reproduire la réponse d'un modèle détaillé de neurone cortical: modèle; potentiels calciques dans les branches des neurones corticaux humains, réalisant le OU exclusif: mesuré sur tranches; apprentissage auto-supervisé: revue des indices & \cite{Payeur2021,Beniaguev2021,Gidon2020,Keller2018} & hypothèse \\
\bottomrule
\end{longtable}}

\section{Chapitre~\ref{sec:sensors}. Les entrées de la machine}

La structure des capteurs est mesurée directement, par des enregistrements du courant de cellules isolées, des vibrations de la cochlée et par le comptage des cellules de la rétine; la limite de sensibilité et la formule du bruit de grenaille découlent de la physique; que les codes des capteurs soient ajustés pour porter le plus d'information est une hypothèse solidement appuyée.

{\footnotesize
\setlength{\tabcolsep}{3pt}\renewcommand{\arraystretch}{1.15}
\begin{longtable}{>{\raggedright}p{0.5cm}>{\raggedright}p{4.0cm}>{\raggedright}p{5.6cm}>{\raggedright}p{2.3cm}>{\raggedright\arraybackslash}p{1.7cm}}
\toprule
N\textsuperscript{o} & Affirmation & Expérience et ce qui est mesuré & Source & Statut \\
\midrule\endfirsthead
\toprule
N\textsuperscript{o} & Affirmation & Expérience et ce qui est mesuré & Source & Statut \\
\midrule\endhead
1 & Le capteur est un transducteur: le récepteur produit un courant, et la sortie des cellules sensorielles de l'œil et de l'oreille est du glutamate vers le bus \nt{GLUT}; les premières cellules ne produisent pas d'impulsions (fig.~\ref{fig:transducer}) & Les bâtonnets et les cônes de la tortue libèrent un acide aminé excitateur, détecté par des canaux glutamatergiques sur un fragment de membrane excisé. Cellules ciliées internes de la cochlée du rat: les courants dans les terminaisons des fibres nerveuses passent par des récepteurs AMPA du glutamate, et la fréquence des libérations croît avec la dépolarisation progressive de la cellule jusqu'à $150$\,Hz & \cite{CopenhagenJahr1989,GlowatzkiFuchs2002} & mesuré \\
2 & Dans l'obscurité, le capteur de lumière répond à un seul photon par un courant d'environ un picoampère & Électrode à succion sur le segment externe d'un bâtonnet de crapaud: les réponses aux éclairs faibles sont quantifiées, événement de $\approx1$\,pA ($5\%$ du courant d'obscurité) avec une dispersion de $0.2$\,pA, efficacité $0.5$. L'homme signale l'arrivée d'un seul photon plus souvent que le hasard & \cite{Baylor1979,Tinsley2016} & mesuré \\
3 & Le capteur de son détecte un déplacement inférieur au nanomètre & Méthode Mössbauer, membrane basilaire de la cochlée du cobaye au site $16$--$19$\,kHz: au seuil de réponse du nerf auditif, la vitesse de la membrane est d'environ $0.04$\,mm/s, soit un déplacement de $\approx0.35$\,nm (notre conversion). On a mesuré la membrane, et non le faisceau du capteur & \cite{Sellick1982,Hudspeth1989} & mesuré \\
4 & Dans l'obscurité, la précision est limitée par le bruit de grenaille des photons~\eqref{eq:photonnoise}; en lumière faible, l'accumulation est plus longue & La formule découle de la statistique de Poisson. Dans le bâtonnet, le bruit en lumière faible coïncide avec la somme d'événements quantiques indépendants; sur un fond lumineux, la réponse à un éclair est plus petite et plus brève. Le chiffre \og dixièmes de seconde\fg{} n'est pas vérifié séparément ici & \cite{Baylor1979} & théorie \\
5 & La plage d'entrée couvre dix ordres de grandeur pour la lumière et douze pour le son, la fréquence des impulsions moins de trois & Pour le son: la réponse de la membrane basilaire à la fréquence caractéristique croît avec une pente descendant jusqu'à $0.2$\,dB/dB dans la bande $40$--$80$\,dB; la compression reste visible au-delà de $100$\,dB. Les nombres d'ordres de grandeur eux-mêmes sont des estimations classiques, sans source dans le chapitre & \cite{Ruggero1997} & mesuré \\
6 & Réponse du capteur~\eqref{eq:nakarushton}; $\sigma$ suit la luminosité moyenne et fait glisser la courbe le long de l'axe logarithmique (fig.~\ref{fig:adaptation}) & La formule a d'abord été ajustée aux potentiels S de la rétine de poisson. Cônes de la tortue: les réponses obéissent à $V=I V_m/(I+\sigma)$ et couvrent $\sim3.5$ ordres de grandeur de luminosité; après $2$--$3$ minutes de fond, la courbe glisse horizontalement en gardant sa largeur. Lien avec la division par l'activité totale: revue & \cite{NakaRushton1966,NormannPerlman1979,Carandini2012} & mesuré \\
7 & Les capteurs transmettent les changements, et leurs codes sont ajustés pour porter le plus d'information par impulsion (proposition~\ref{prop:sensors}) & Le principe a été proposé en théorie. L'indice le plus fort: chez la mouche, la caractéristique \og contraste $\to$ réponse\fg{} des neurones de la première couche coïncide avec la distribution cumulée des contrastes des scènes naturelles, ce qui maximise l'information & \cite{Barlow1961,Laughlin1981} & hypothèse \\
8 & Les capteurs ON et OFF forment un code à deux fils; la caméra événementielle le reproduit dans le silicium & Revue d'expériences: les voies ON et OFF de la rétine se séparent dès la première synapse et peuvent être supprimées séparément. Caméra: $128\times128$ pixels sans images, plage de $120$\,dB, latence de $15$\,\textmu s & \cite{Schiller1992,Lichtsteiner2008,Gallego2022} & mesuré \\
9 & L'œil compte $\sim10^8$ capteurs de lumière et $\sim10^6$ neurones de sortie: compression d'un facteur $\sim100$ & Comptage sur des rétines humaines entières: en moyenne $92$ millions de bâtonnets et $4.6$ millions de cônes; de $0.7$ à $1.5$ million de cellules de sortie (ganglionnaires) selon les yeux & \cite{Curcio1990,CurcioAllen1990} & mesuré \\
10 & Chez la souris, l'œil a plus de trente canaux de sortie & Souris: imagerie calcique biphotonique de plus de $11\,000$ cellules de sortie et regroupement des réponses, nettement plus de $30$ types fonctionnels. Chez l'homme, le nombre n'est pas mesuré & \cite{Baden2016} & mesuré \\
11 & L'œil du cobaye transmet $\sim875\,000$ bit/s; rapporté à l'œil humain, cela donne $\sim10^7$ bit/s & Cobaye, matrice multiélectrode, mouvement dans des scènes naturelles: $6$--$13$ bit/s par cellule, $\sim875\,000$ bit/s pour $\sim10^5$ cellules de sortie. Pour l'homme ($\sim10^6$ cellules), $\sim10^7$ bit/s: extrapolation des auteurs & \cite{Koch2006} & mesuré \\
12 & L'oreille est un banc de filtres passe-bande avec une carte des fréquences presque logarithmique (fig.~\ref{fig:filterbank}) & Le lieu de vibration maximale de la membrane basilaire dépend de la fréquence; la fonction \og fréquence--lieu\fg{} est presque exponentielle et décrit sous une seule forme les cochlées de l'homme et d'animaux vivants & \cite{Greenwood1990,Robles2001} & mesuré \\
13 & Les résonateurs sont actifs: une partie des capteurs amplifie les vibrations faibles d'un facteur de plusieurs milliers en amplitude ($66$--$76$\,dB); le gain diminue quand le volume augmente & Vibrométrie laser à la base de la cochlée du chinchilla: pour des sons faibles à la fréquence caractéristique, gain de $66$--$76$\,dB par rapport à l'étrier; la mort réduit la sensibilité de $60$--$81$\,dB ($10^3$--$10^4$ fois en amplitude) et supprime la compression. La source de puissance est constituée par les cellules ciliées externes & \cite{Ruggero1997,Robles2001} & mesuré \\
14 & L'amplificateur de l'oreille est placé au bord de l'auto-oscillation & Calcul: près d'une bifurcation de Hopf, compression de la plage, accord fin et sons de combinaison sont inévitables, soit toutes les non-linéarités connues de l'audition. Que la cochlée soit réglée exactement ainsi n'est pas démontré directement & \cite{Eguiluz2000} & hypothèse \\
15 & Aux basses fréquences, les impulsions sont en phase avec le son (chez le cobaye, le verrouillage faiblit à partir de $\sim600$\,Hz et reste perceptible jusqu'à $\sim3.5$\,kHz), et la direction s'en déduit; aux hautes fréquences, il reste le code de position & Nerf auditif du cobaye: le verrouillage de phase commence à décliner vers $600$\,Hz et disparaît au-dessus de $3.5$\,kHz; chez le chat et le singe, la limite est environ une octave plus haut. Différence de temps d'arrivée aux deux oreilles: revue & \cite{PalmerRussell1986,Grothe2010} & mesuré \\
16 & Autres capteurs (tabl.~\ref{tab:sensors}): l'odorat a $\sim400$ types de récepteurs, un par neurone, et un code combinatoire; le goût passe en partie par l'ATP et la sérotonine; le toucher a des capteurs à adaptation rapide et lente; le chaud et le froid sont séparés & Souris: un récepteur reconnaît de nombreuses odeurs, une odeur active de nombreux récepteurs, des odeurs différentes activent des ensembles différents. Sans récepteurs de l'ATP, le nerf gustatif ne répond pas; les bourgeons du goût libèrent de la sérotonine. Enregistrement de fibres cutanées isolées de la main humaine: quatre types selon la vitesse d'adaptation. Le canal du froid est distinct de ceux du chaud & \cite{Mombaerts2004,Malnic1999,Finger2005,Huang2005,JohanssonVallbo1979,McKemy2002} & mesuré \\
\bottomrule
\end{longtable}}

\section{Chapitre~\ref{sec:recognition}. Reconnaissance}

La vitesse de reconnaissance, la structure des premiers étages, la lecture linéaire de la réponse et l'apprentissage sur la continuité du temps sont mesurés chez l'homme et le singe; la réduction de toute la chaîne à deux opérations et l'apprentissage sur la vidéo sont vérifiés sur les modèles du livre; les explications des illusions vont du bien fondé au discuté.

{\footnotesize
\setlength{\tabcolsep}{3pt}\renewcommand{\arraystretch}{1.15}
\begin{longtable}{>{\raggedright}p{0.5cm}>{\raggedright}p{4.0cm}>{\raggedright}p{5.6cm}>{\raggedright}p{2.3cm}>{\raggedright\arraybackslash}p{1.7cm}}
\toprule
N\textsuperscript{o} & Affirmation & Expérience et ce qui est mesuré & Source & Statut \\
\midrule\endfirsthead
\toprule
N\textsuperscript{o} & Affirmation & Expérience et ce qui est mesuré & Source & Statut \\
\midrule\endhead
1 & Sous translation, la comparaison pixel par pixel avec un modèle ne fait pas mieux que pile ou face; la paire d'étages~\eqref{eq:scstage} reconnaît la figure à n'importe quelle position & \texttt{models/\allowbreak recognition\_models.py}, \og rétine\fg{} de $24$ points, $4000$ essais: modèle, $0.49$, $0.51$, $0.52$ pour une proportion de points inversés de $0$, $0.1$, $0.2$; paire $S$ et $C$, $1.00$, $0.86$, $0.70$ & déduction du chapitre & modèle \\
2 & Un animal sur une image est reconnu en $\sim150$\,ms, en un seul passage à travers une dizaine d'étages de $10$--$20$\,ms & Homme, tâche \og animal ou pas\fg{} sur des photos présentées $20$\,ms: les potentiels évoqués des deux réponses divergent après $\sim150$\,ms. Singe: le temps de première réponse croît d'étage en étage (V1 d'abord, puis V2 et V4). Le nombre d'étages et \og zéro ou une impulsion par étage\fg{} sont déduits de ces chiffres & \cite{Thorpe1996,Schmolesky1998} & mesuré \\
3 & L'étage $S$ est un ET sur les parties, l'étage $C$ un maximum sur les positions (proposition~\ref{prop:conveyor}) & Chat, cortex visuel primaire: les cellules simples répondent à un bord d'orientation donnée à un endroit donné, les cellules complexes à la même orientation sur une zone plus grande. Enregistrement intracellulaire de cellules complexes: la réponse à deux barres est, pour beaucoup de cellules, proche de la plus grande des deux réponses, en moyenne très proche de l'opération max. Maximum par inhibition mutuelle: proposition~\ref{prop:wta}; division par l'activité totale: revue & \cite{HubelWiesel1962,Lampl2004,Carandini2012} & mesuré \\
4 & Plus haut dans la chaîne, les combinaisons sont plus grandes et plus complexes, et la réponse dépend de moins en moins de la position & Singe, simplification des stimuli jusqu'au \og trait critique\fg{}: les cellules de V4 et du cortex inférotemporal postérieur répondent à des combinaisons complexes de forme, de couleur et de texture avec un petit champ; dans le cortex inférotemporal antérieur, le champ est plus grand. L'alternance de $S$ et $C$ dans le modèle reproduit ce tableau & \cite{KobatakeTanaka1994,Riesenhuber1999} & mesuré \\
5 & Les réseaux convolutifs reproduisent cette alternance et prédisent le mieux les réponses des étages supérieurs; l'objet se lit linéairement à partir des fréquences d'un groupe de neurones & Réseau entraîné à reconnaître, sans ajustement au cerveau: la couche supérieure prédit les réponses des neurones du cortex inférotemporal du singe, les couches intermédiaires celles de V4. Un classifieur linéaire sur l'activité de $\sim100$ cellules prises au hasard, dans des fenêtres dès $12.5$\,ms, reconnaît l'objet et sa catégorie à différentes positions et tailles & \cite{Fukushima1980,Yamins2014,Hung2005} & mesuré \\
6 & La règle de Hebb avec trace~\eqref{eq:traceinvariance} apprend l'invariance sur une vidéo sans étiquettes si la trace dure au moins $\sim20$\,ms & L'idée des traits lents et de la règle avec trace relève de la théorie. \texttt{models/\allowbreak recognition\_models.py}, deux neurones $C$, $16$ entrées, $10$ essais, pause entre objets de $0.3$\,s. Pour $\tau_{\text{tr}}=5$\,ms, coïncidence avec la figure de $0.52$ en moyenne, pour $10$\,ms, $0.56$; pour $20$, $50$ et $200$\,ms, $1.00$ dans tous les essais (à $30$\,ms, un essai sur dix se trompe) & \cite{Foldiak1991,WiskottSejnowski2002} & modèle \\
7 & Dans le cerveau, l'invariance s'apprend sur la continuité du temps & Singe: on substituait discrètement l'objet pendant un saut de l'œil vers un endroit donné du champ; l'invariance de position des neurones du cortex inférotemporal se modifiait conformément à la substitution, de façon significative dès une heure. Qu'il s'agisse de la règle principale de l'apprentissage visuel est une hypothèse & \cite{LiDiCarlo2008} & mesuré \\
8 & Mouvement: détecteurs de direction, puis la même paire ET/OU sur de grandes zones & Détecteurs de direction dans la rétine du lapin; dans l'aire MT du singe, les $110$ neurones enregistrés sont tous sélectifs à la direction, et la réponse au mouvement y est plus forte que dans V1 & \cite{Barlow1965,Albright1984} & mesuré \\
9 & Les étages supérieurs envoient une prédiction vers le bas, l'erreur monte & Le modèle explique les effets extra-classiques des champs du cortex visuel primaire; certaines observations concordent (revue), mais ce n'est pas démontré comme principe général du pipeline & \cite{RaoBallard1999,Keller2018} & hypothèse \\
10 & Les images difficiles exigent des rétroactions et plusieurs tours & Singe: pour les images \og difficiles\fg{}, la décision dans le cortex inférotemporal se forme $\sim30$\,ms plus tard; ces réponses tardives sont mal prédites par un réseau direct, mieux par des réseaux récurrents ou plus profonds & \cite{Kar2019} & mesuré \\
11 & Une falsification invisible des pixels trompe le réseau; la vision humaine est bien plus robuste, mais pas invulnérable & Sur les réseaux: une petite perturbation choisie tout exprès change la réponse; l'exemple \og panda $\to$ gibbon\fg{} vient du deuxième article. Chez l'homme, en présentation brève, des falsifications qui se transfèrent bien d'un réseau à l'autre infléchissent le choix & \cite{Szegedy2014,Goodfellow2015,Elsayed2018} & mesuré \\
12 & Illusions d'attente: une entrée peu fiable cède à l'attente; ce qui est pâle bouge \og plus lentement\fg{}, la robe est vue de façons différentes (section~\ref{sec:illusions}) & Des réseaux de barres moins contrastés semblent bouger plus lentement. Enquête auprès de $1401$ personnes: la robe est vue selon trois variantes stables, un indice sur l'éclairage fait basculer la couleur; chez $\sim13\,000$ observateurs, c'est l'hypothèse sur l'éclairage qui décide. Un modèle bayésien avec attente de mouvements lents explique plusieurs illusions de mouvement & \cite{Thompson1982,LaferSousa2015,Wallisch2017,Weiss2002,Purves2011} & hypothèse \\
13 & Réglage du gain: cascade, inclinaison et contraste; la grille d'Hermann ne s'explique pas par l'inhibition des voisins & Les détecteurs de direction de la rétine du lapin, lors d'un mouvement prolongé, baissent leur fréquence et, après l'arrêt, descendent sous le niveau de fond. L'illusion d'inclinaison est reproduite par un modèle de division par l'activité des voisins. Des modifications de la grille qui ne changent pas la réponse des neurones de sortie de la rétine suppriment les taches: l'explication par inhibition des voisins est rejetée & \cite{BarlowHill1963,Schwartz2009,SchillerCarvey2005} & mesuré \\
14 & Compensation des délais: un objet en mouvement semble en avance sur l'éclair & L'illusion est mesurée. La psychophysique contredit à la fois le prolongement du mouvement et la différence de délais; une explication a posteriori est proposée, par les événements survenant $\sim80$\,ms après l'éclair. Le débat n'est pas tranché & \cite{Nijhawan1994,EaglemanSejnowski2000} & hypothèse \\
15 & Des \og serpents\fg{} immobiles tournent: la différence de délais~\eqref{eq:latency} est lue par les détecteurs de direction & L'illusion fonctionne aussi sans couleur; des paires d'éléments voisins la produisent isolément; des neurones corticaux du singe sélectifs à la direction donnent une réponse orientée à ces paires immobiles. Chez l'homme, la rotation est déclenchée par les microsaccades et les clignements & \cite{Conway2005,OteroMillan2012} & mesuré \\
16 & Images ambiguës: inhibition mutuelle, fatigue et bruit; les bascules sont surtout dues au bruit & Modèle de groupes de neurones en compétition: les bascules y sont causées par le \emph{bruit}, sans bruit il n'y en a pas; il explique la hausse de la fréquence des bascules avec l'intensité du stimulus. La fatigue y joue un rôle moindre & \cite{MorenoBote2007} & hypothèse \\
17 & Une partie des illusions découle de la tâche qu'apprend la machine & Un réseau entraîné à prédire l'image suivante d'une vidéo \og voit\fg{} la rotation d'anneaux immobiles et ne la voit pas sur les images témoins; des réseaux de débruitage et d'augmentation de la netteté cèdent aux illusions de couleur et de luminosité, comme l'homme & \cite{Watanabe2018,GomezVilla2020} & mesuré \\
\bottomrule
\end{longtable}}

\section{Chapitre~\ref{sec:muscles}. La sortie vers les muscles}

La structure de la sortie --- unités motrices, marge de sécurité de la synapse, activation par taille, paire de muscles, boucle d'étirement et générateurs de rythme --- est mesurée directement; la condition de stabilité de la boucle à retard est un théorème; le modèle interne du corps est une hypothèse bien fondée; les chiffres des figures sont calculés avec le modèle du livre.

{\footnotesize
\setlength{\tabcolsep}{3pt}\renewcommand{\arraystretch}{1.15}
\begin{longtable}{>{\raggedright}p{0.5cm}>{\raggedright}p{4.0cm}>{\raggedright}p{5.6cm}>{\raggedright}p{2.3cm}>{\raggedright\arraybackslash}p{1.7cm}}
\toprule
N\textsuperscript{o} & Affirmation & Expérience et ce qui est mesuré & Source & Statut \\
\midrule\endfirsthead
\toprule
N\textsuperscript{o} & Affirmation & Expérience et ce qui est mesuré & Source & Statut \\
\midrule\endhead
1 & Unité motrice: un neurone \nt{ACET} commande un groupe de fibres, de quelques centaines à deux mille environ; dans les muscles à réglage fin, les unités sont plus petites & Muscles humains, comptage des axones moteurs et des fibres musculaires: en moyenne environ $340$ fibres par neurone dans un muscle interosseux de la main, environ $410$ dans le muscle brachio-radial, environ $1930$ dans le gastrocnémien médial. Le chapitre ne donne pas de chiffre précis pour les plus petites unités. & \cite{Feinstein1955mu} & mesuré \\
2 & La synapse sur le muscle libère plusieurs fois plus de neurotransmetteur qu'il n'en faut pour déclencher la fibre & Revue des mesures sur les synapses neuromusculaires: la marge de sécurité (rapport du neurotransmetteur libéré au seuil) vaut en général $3$--$5$ chez les mammifères adultes; chez l'homme, elle repose davantage sur les replis de la membrane du destinataire que sur l'importance de la libération. & \cite{WoodSlater2001} & mesuré \\
3 & La secousse~\eqref{eq:twitch} monte en un temps $T_c$ de l'ordre de $50\ms$ & Unités motrices isolées de la main humaine lors d'un effort volontaire, force moyennée sur les impulsions de l'unité: temps de montée de $30$--$100\ms$, inférieur à $70\ms$ pour $80\,\%$ des unités. La fonction $(t/T_c)e^{1-t/T_c}$ est une approximation tirée d'un modèle de pool d'unités. & \cite{MilnerBrown1973rec, Fuglevand1993} & mesuré \\
4 & Somme des secousses~\eqref{eq:forcesum}: à $5$, $15$ et $40$ imp./s, force de $0.2$--$1.1F_1$, $1.8$--$2.2F_1$ et environ $5.4F_1$ (fig.~\ref{fig:tetanus}) & Calcul avec \texttt{models/\allowbreak muscle\_\allowbreak models.py}, $T_c=50\ms$: $0.21$--$1.10$, $1.76$--$2.19$ et $5.28$--$5.49\,F_1$ sur les dernières $0.3\,\text{s}$. L'addition linéaire est une simplification: le modèle ne contient pas la saturation d'un vrai muscle. & \texttt{models/\allowbreak muscle\_\allowbreak models.py} & modèle \\
5 & Les unités s'activent par ordre de taille; les plus faibles sont cent fois plus faibles que les plus fortes & Racines ventrales du chat: sous une entrée réflexe croissante, les premiers neurones à décharger sont ceux dont les impulsions sont petites dans la racine, c'est-à-dire les petits. Muscle interosseux de la main humaine: force de secousse des unités de $0.1$ à $10$~g, croissant presque linéairement avec la force à laquelle l'unité s'active. & \cite{Henneman1957, MilnerBrown1973rec} & mesuré \\
6 & Le pas de force est proportionnel à la force, $\Delta F/F\approx q-1$~\eqref{eq:sizeprinciple}: virgule flottante & Déduit dans le chapitre de la progression géométrique des forces. Concorde avec l'expérience: aux efforts élevés, bien moins d'unités nouvelles s'activent pour un même incrément de force. & déduction du chapitre; \cite{MilnerBrown1973rec} & théorie \\
7 & La dispersion de la force croît proportionnellement à la force & Effort isométrique volontaire d'un muscle du pouce: l'écart type de la force croît linéairement avec la force moyenne; pour le même effort provoqué par stimulation électrique, la dispersion est constante. Modèle de pool: la croissance linéaire vient de l'activation des unités par ordre de force. & \cite{Jones2002sdn} & mesuré \\
8 & La fluidité des mouvements découle d'un bruit dépendant du signal & Modèle: la trajectoire qui minimise la dispersion du point d'arrivée sous un tel bruit reproduit les trajectoires mesurées de l'œil et du bras. Aucune expérience directe ne distingue cette cause des autres. & \cite{HarrisWolpert1998} & hypothèse \\
9 & La différence des forces d'une paire de muscles fixe le couple, leur somme la raideur~\eqref{eq:antagonists} & Théorie du contrôle de l'impédance par co-contraction. Bras humain: la raideur de chaque articulation, mesurée par de petites poussées, est à peu près proportionnelle au couple qui s'y exerce; différents rapports de co-contraction changent la forme et l'orientation de l'ellipse de raideur de la main. & \cite{Hogan1984, GomiOsu1998stiff} & mesuré \\
10 & Le capteur d'étirement excite son propre neurone \nt{ACET} et, par un intermédiaire inhibiteur, coupe l'antagoniste (fig.~\ref{fig:antagonists}) & Moelle épinière du chat: un seul intermédiaire, excité par les fibres des capteurs d'étirement, produit des potentiels inhibiteurs monosynaptiques dans les neurones \nt{ACET} du muscle antagoniste, délai synaptique $0.28$--$0.42\ms$. Le neurotransmetteur de l'intermédiaire (la glycine) n'a pas été déterminé dans cette expérience. & \cite{Jankowska1972Ia} & mesuré \\
11 & Retard de la boucle: $25$--$30\ms$ pour la boucle spinale, une fois et demie à trois fois plus par le cortex, $100$--$200\ms$ par l'œil & Bras humain, poussée sur l'articulation: réponse musculaire courte après $20$--$45\ms$, longue après $45$--$105\ms$, volontaire après $120$--$180\ms$. Décalage de la position visible du doigt pendant le mouvement: correction après $160\ms$ en moyenne. & \cite{Pruszynski2008rapid, SaundersKnill2003vis} & mesuré \\
12 & $\dd x/\dd t=-Gx(t-d)$ est stable pour $Gd<\pi/2$~\eqref{eq:delaystable}; à la limite, la fréquence est $1/(4d)$ & Théorème sur les racines d'une équation à retard. Vérification avec \texttt{models/\allowbreak muscle\_\allowbreak models.py}, $d=30\ms$: à $3\,\text{s}$, amplitude de $3\cdot10^{-6}$ pour $G=40$, $0.13$ pour $G=50$, $38$ pour $G=55$ par seconde (limite $52$). & \cite{Hayes1950} & théorie \\
13 & Tremblement physiologique de $8$--$12$~Hz; la boucle d'étirement est l'une de ses sources & Revue des mesures: un petit tremblement vers $10$~Hz existe chez les personnes en bonne santé; son origine est mixte: oscillations centrales, propriétés des décharges des unités, résonance mécanique et résonance de la boucle réflexe, et la part dominante varie selon les conditions. & \cite{McAuleyMarsden2000} & hypothèse \\
14 & Le cerveau prédit le résultat d'une commande grâce à un modèle interne du corps & Un sujet tient une charge entre deux doigts et bouge le bras: sous charge inertielle, visqueuse et élastique, la force de préhension varie en même temps que la charge, et non avec le retard de la boucle; la charge est donc prédite. Une revue rassemble ces expériences; on n'observe pas directement le modèle lui-même dans les neurones. & \cite{FlanaganWing1997grip, WolpertGhahramani2000} & hypothèse \\
15 & Le rythme de la marche, de la respiration, de la mastication est produit par des réseaux locaux sous entrée constante & Revue d'expériences: un système nerveux isolé (moelle épinière, ganglions) produit un motif moteur rythmique sans entrée rythmique et sans rétroaction des muscles. & \cite{MarderBucher2001} & mesuré \\
16 & L'inhibition mutuelle avec fatigue~\eqref{eq:halfcentre} donne un rythme de période $0.84\,\text{s}\approx2\tau_a$ (fig.~\ref{fig:halfcentre}) & Théorème: un réseau à inhibition mutuelle et adaptation sans état stable oscille nécessairement sous entrée constante. Calcul avec \texttt{models/\allowbreak muscle\_\allowbreak models.py} ($I=1$, $\beta=2$, $b=2$, $\tau=20\ms$, $\tau_a=0.4\,\text{s}$): période de $0.84\,\text{s}$. Que les vrais générateurs soient construits ainsi n'est pas démontré. & \cite{Matsuoka1985osc} & modèle \\
\bottomrule
\end{longtable}}

\section{Chapitre~\ref{sec:pain}. La douleur}

Les capteurs de la douleur, les deux fils, le portillon de la moelle épinière, le bouton de volume descendant, la sensibilisation et la capture de l'attention sont mesurés directement; les formules du portillon et de la sensibilisation sont des modèles simplifiés du livre; la règle selon laquelle la machine choisit le volume de l'alarme est une hypothèse.

{\footnotesize
\setlength{\tabcolsep}{3pt}\renewcommand{\arraystretch}{1.15}
\begin{longtable}{>{\raggedright}p{0.5cm}>{\raggedright}p{4.0cm}>{\raggedright}p{5.6cm}>{\raggedright}p{2.3cm}>{\raggedright\arraybackslash}p{1.7cm}}
\toprule
N\textsuperscript{o} & Affirmation & Expérience et ce qui est mesuré & Source & Statut \\
\midrule\endfirsthead
\toprule
N\textsuperscript{o} & Affirmation & Expérience et ce qui est mesuré & Source & Statut \\
\midrule\endhead
1 & Seuil élevé: les seuils d'échauffement des différents capteurs de la douleur vont de $41^\circ$ à $46$--$48^\circ$ & Enregistrement de fibres isolées. Peau de singe: seuil médian de $41^\circ$ pour les capteurs lents (C), $46^\circ$ pour les capteurs rapides de type II, plus de $53^\circ$ pour le type I. Peau humaine: $40.7^\circ$ pour les capteurs C qui répondent aussi à la pression, $48^\circ$ pour ceux qui n'y répondent pas. & \cite{Treede1995heat, Weidner1999cfib} & mesuré \\
2 & Les capteurs de la douleur s'adaptent bien moins que les capteurs du toucher et deviennent plus sensibles après une lésion légère; un affaiblissement après un fort échauffement est mesuré pour un seul type & Brûlure légère de la peau ($50^\circ$, $100\,\text{s}$): après $5$--$10$~min, le seuil de douleur chez l'homme et le seuil des capteurs C chez le singe sont abaissés de $2$--$6^\circ$; les autres capteurs cutanés ne présentent pas ce décalage. Chez les capteurs rapides de type II, la réponse est supprimée après un fort échauffement. La comparaison de l'accoutumance avec les capteurs du toucher est une estimation générale, non étayée ici par une expérience propre. & \cite{LaMotte1982hyper, Treede1995heat} & mesuré \\
3 & Deux fils: rapide, $5$--$30$~m/s; lent, $0.5$--$2$~m/s; temps d'arrivée $t=L/v$~\eqref{eq:paintime} & Vitesse de conduction: capteurs rapides de la douleur du singe, en moyenne $15$ et $25$~m/s (deux types); capteurs C humains, $0.8$--$1.0$~m/s. Pour $L=1$~m, cela fait quelques dizaines de millisecondes et environ une seconde. & \cite{Treede1995heat, Weidner1999cfib} & mesuré \\
4 & La douleur arrive deux fois: d'abord rapide et précise, puis sourde et durable & Temps de réaction à l'échauffement de l'avant-bras humain: de $1100$ à $400\ms$ quand la température monte; un fort échauffement de la peau velue donne une double douleur, celui de la paume non. Seules des fibres plus rapides que $6$~m/s peuvent porter la première douleur; par la latence, elle correspond aux capteurs rapides de type II, absents de la paume. Le partage des rôles (\og où\fg{} et \og à quel point c'est grave\fg{}) est une interprétation. & \cite{CampbellLaMotte1983first, Treede1995heat} & mesuré \\
5 & Portillon: le neurone de sortie de la moelle reçoit la douleur et le toucher, l'intermédiaire inhibiteur est excité par le toucher (fig.~\ref{fig:gate}) & Le schéma du portillon a été proposé comme théorie. Souris, marquage génétique et suppression de groupes de neurones: les neurones excitateurs à somatostatine sont nécessaires à la douleur mécanique; les neurones inhibiteurs à dynorphine empêchent l'entrée des capteurs du toucher de les atteindre; sans ces neurones, un toucher léger provoque de la douleur. & \cite{MelzackWall1965, Duan2014} & mesuré \\
6 & Le toucher atténue la douleur & Homme, impulsions laser douloureuses sur la main: un toucher simultané réduit la douleur ressentie et la capacité à distinguer l'énergie des impulsions; l'effet faiblit avec la distance entre le point de toucher et le point de douleur au sein d'un même segment cutané. & \cite{Mancini2014touch} & mesuré \\
7 & Hypothèse de la théorie du portillon: une douleur forte éteint elle-même l'intermédiaire inhibiteur, et le portillon s'ouvre grand & Présenté dans le chapitre comme une hypothèse de la théorie originale. Le travail sur la souris décrit comment le portillon empêche le toucher d'entrer dans le canal de la douleur; l'extinction de l'intermédiaire par la douleur n'y est pas montrée. & \cite{MelzackWall1965} & hypothèse \\
8 & Portillon~\eqref{eq:gate}: seuil de $0.18$ sans toucher, $0.86$ avec toucher, $0.11$ pour $g=2$; pour $\nu=1$, le toucher réduit la sortie de $1$ à $0.25$, pour $\nu=2$ il est sans effet; le toucher seul ne passe pas (fig.~\ref{fig:gatecurves}) & Le calcul avec \texttt{models/\allowbreak pain\_\allowbreak gate.py} et les poids du texte reproduit tous ces nombres. & \texttt{models/\allowbreak pain\_\allowbreak gate.py} & modèle \\
9 & Les voies descendantes du tronc cérébral modifient le gain du portillon dans les deux sens & Rat, bulbe rachidien, enregistrement pendant l'échauffement de la queue: certains neurones déchargent avant le retrait (cellules \og ON\fg{}), d'autres se taisent (cellules \og OFF\fg{}); une faible stimulation de cette région supprime le réflexe. Revue: les mêmes circuits facilitent et inhibent la transmission de la douleur, et les opioïdes agissent par eux. & \cite{Fields1983rvm, Fields2004} & mesuré \\
10 & L'attente d'un soulagement atténue la douleur par le canal opioïde & Patients souffrant d'une douleur aiguë: chez ceux dont l'attente d'un soulagement avait réduit la douleur, le blocage des récepteurs opioïdes a ramené la douleur au niveau des autres; chez les autres, le blocage ne changeait pas la douleur. & \cite{Levine1978} & mesuré \\
11 & Le cerveau choisit le volume de l'alarme selon l'intérêt de l'interruption & Il n'existe pas d'expérience mesurant la règle du choix du volume. & --- & hypothèse \\
12 & Sensibilisation: après une lésion, la sortie du portillon augmente et le seuil baisse, en partie grâce à la moelle épinière; elle persiste après l'arrêt du stimulus nocif & Rat: après une lésion cutanée, le seuil du réflexe de flexion baisse et la réponse augmente; l'analyse électrophysiologique montre qu'une partie de cette hausse naît dans la moelle elle-même. Un stimulus nocif provoque des troubles durables de la sensibilité qui survivent au stimulus. Le chapitre ne donne pas le temps de décroissance exact. & \cite{Woolf1983} & mesuré \\
13 & La sensibilisation~\eqref{eq:sensitization} est une réaction positive; si la boucle est forte, l'alarme peut se verrouiller après la guérison & Découle du modèle~\eqref{eq:sensitization} (exercice en fin de chapitre). Aucune expérience ne montre qu'une douleur persistante après guérison soit une boucle verrouillée de ce type précis. & déduction du chapitre & hypothèse \\
14 & La douleur est une récompense négative, et l'apprentissage par la douleur suit l'erreur de différence temporelle & IRM, sujets humains, chaînes de signes annonciateurs de la douleur: l'activité du striatum ventral et de l'insula antérieure coïncide avec les signaux d'apprentissage par différence temporelle. Singe: une partie des neurones \nt{DOPA} est excitée par un souffle d'air désagréable et son signe annonciateur, une autre partie est inhibée. & \cite{Seymour2004, Matsumoto2009} & mesuré \\
15 & La douleur interrompt le travail en cours & Sujets humains, stimulus électrique douloureux et sondes sonores: la réponse à une sonde présentée $100\ms$ après le début de la douleur est nettement ralentie; à $1.5\,\text{s}$, il n'y a plus de différence avec le contrôle. Une revue rassemble ces expériences dans un modèle de capture de l'attention. & \cite{Crombez1996disrupt, EcclestonCrombez1999} & mesuré \\
16 & Le retrait se referme dans la moelle épinière & Rat: les neurones des couches profondes de la corne dorsale reproduisent le motif spatial du réflexe de retrait de muscles isolés; on ne parvient pas à les exciter de façon antidromique depuis la région cervicale: ce sont donc des intermédiaires spinaux. & \cite{Schouenborg1995wr} & mesuré \\
\bottomrule
\end{longtable}}

\section{Chapitre~\ref{sec:motorlearning}. L'apprentissage moteur}

La règle des moindres carrés et l'avantage d'un fil d'erreur séparé sont des théorèmes; la structure du circuit à fil d'erreur, l'affaiblissement par coïncidence, l'ajustement de la trace au retard, l'effet consécutif et le retour spontané de l'habileté sont mesurés; que le circuit entier fonctionne comme un apprentissage supervisé et construise un modèle interne est l'hypothèse dominante; les chiffres du modèle à deux processus sont calculés avec le modèle du livre.

{\footnotesize
\setlength{\tabcolsep}{3pt}\renewcommand{\arraystretch}{1.15}
\begin{longtable}{>{\raggedright}p{0.5cm}>{\raggedright}p{4.0cm}>{\raggedright}p{5.6cm}>{\raggedright}p{2.3cm}>{\raggedright\arraybackslash}p{1.7cm}}
\toprule
N\textsuperscript{o} & Affirmation & Expérience et ce qui est mesuré & Source & Statut \\
\midrule\endfirsthead
\toprule
N\textsuperscript{o} & Affirmation & Expérience et ce qui est mesuré & Source & Statut \\
\midrule\endhead
1 & La règle~\eqref{eq:lms} suit en moyenne le gradient du carré de l'erreur & Résultat de la théorie des filtres adaptatifs: une correction proportionnelle à l'entrée et à l'erreur de sortie est une descente stochastique du gradient de l'erreur quadratique moyenne. & \cite{WidrowHoff1960} & théorie \\
2 & Avec un fil d'erreur par sortie, la vitesse ne dépend pas du nombre de sorties; avec un seul nombre commun, elle baisse avec le nombre de neurones bruités (proposition~\ref{prop:teachers}) & Courbes d'apprentissage d'un réseau linéaire: l'apprentissage par perturbations aléatoires des neurones est plus lent que le gradient exact d'un facteur égal au nombre de neurones perturbés. & \cite{Werfel2005} & théorie \\
3 & Le circuit à fil d'erreur fonctionne comme un apprentissage supervisé (proposition~\ref{prop:errorwire}) & Théories déduites de la structure du circuit. Les expériences des lignes 7--10 concordent avec elles; que le circuit entier fonctionne exactement ainsi n'est pas démontré. & \cite{Marr1969, Albus1971} & hypothèse \\
4 & Couche d'expansion: environ $7\cdot10^{10}$ petits neurones \nt{GLUT}, plus que dans tout le reste du cerveau & Comptage des noyaux dans une suspension de tissu dissous: $69\cdot10^9$ neurones dans le cervelet humain sur $86\cdot10^9$ dans tout le cerveau. Stéréologie: $101\cdot10^9$ petits neurones dans le cervelet d'hommes âgés. & \cite{Azevedo2009, Andersen1992cereb} & mesuré \\
5 & Augmenter la dimension de l'entrée rend linéaire la dépendance voulue & Théorème sur le nombre de dichotomies: une somme pondérée de $n$ entrées avec seuil réalise un étiquetage arbitraire d'environ $2n$ vecteurs en position générale. Que le cervelet exploite précisément cela fait partie de l'hypothèse de la ligne~3. & \cite{Cover1965} & théorie \\
6 & Environ $3\cdot10^7$ neurones \nt{GABA} de sortie, chacun avec de l'ordre de $10^5$ entrées & Stéréologie du cervelet humain: $30.5\cdot10^6$ neurones de sortie. Microscopie électronique, rat: environ $1.75\cdot10^5$ entrées de la couche d'expansion par neurone de sortie. Le neurotransmetteur inhibiteur des neurones de sortie: dans la revue. & \cite{Andersen1992cereb, NapperHarvey1988, Ito2001} & mesuré \\
7 & Exactement un fil d'erreur par neurone de sortie, environ une fois par seconde, avec chaque fois une décharge puissante & Revue d'enregistrements: chaque neurone de sortie reçoit une seule entrée \nt{GLUT} grimpante, qui provoque une impulsion complexe à une fréquence de l'ordre de $1$~Hz. & \cite{Ito2001} & mesuré \\
8 & La probabilité de la décharge dépend de la direction de l'erreur de mouvement & Singe, région oculomotrice du cervelet, adaptation des saccades: la direction de l'erreur visuelle fixe la probabilité de l'impulsion complexe, sa grandeur fixe son instant; la décharge modifie les impulsions simples à l'essai suivant et déplace le mouvement le long de l'erreur préférée de la cellule. & \cite{Herzfeld2018} & mesuré \\
9 & La coïncidence de la décharge d'erreur avec une entrée active affaiblit cette entrée ($-\eta e_i x_j$) & Revue d'expériences: la stimulation conjointe d'une entrée de la couche d'expansion et du fil d'erreur provoque un affaiblissement durable de cette entrée; la stimulation de l'entrée seule, non. & \cite{Ito2001} & mesuré \\
10 & La durée de la trace est ajustée au retard de l'erreur: pour un retard de $100\ms$, c'est l'entrée active $100\ms$ avant qui s'affaiblit le plus & Tranches et souris vivantes: dans la région chargée de l'apprentissage oculomoteur, l'affaiblissement est étroitement accordé sur un intervalle d'environ $120\ms$ entre l'entrée et l'erreur, soit le temps que met le signal d'erreur dans cette tâche; dans une autre région, différentes cellules sont accordées sur différents intervalles. & \cite{Suvrathan2016} & mesuré \\
11 & Le circuit est un filtre adaptatif~\eqref{eq:adaptivefilter} qui apprend un modèle interne du corps & Modèle déduit de la structure du circuit. Aucune expérience directe n'a mesuré le filtre appris comme fonction de transfert du corps. & \cite{Fujita1982} & hypothèse \\
12 & La prédiction soustrait les conséquences de nos propres mouvements, c'est pourquoi on ne peut pas se chatouiller soi-même & Sujets humains: un toucher par soi-même semble moins chatouilleux que le même toucher par autrui; en IRM, le cortex somatosensoriel répond plus faiblement au toucher par soi-même, et le cervelet est moins actif quand le mouvement touche la peau. L'explication par soustraction de la prédiction est une hypothèse. & \cite{Blakemore1998} & hypothèse \\
13 & Sous une force extérieure, le raté diminue, et après sa suppression la main rate de l'autre côté & Des sujets tendent le bras en tenant la poignée d'un robot dans un champ de forces: les trajectoires redeviennent droites; après la suppression du champ, l'effet consécutif est presque l'image en miroir des premières déformations; il se transfère aussi à des régions non entraînées de l'espace. & \cite{ShadmehrMussaIvaldi1994} & mesuré \\
14 & Deux processus apprennent~\eqref{eq:tworate}; après une perturbation inverse, l'habileté revient d'elle-même & Sujets humains, champ de forces, puis bref champ inverse et essais à erreur bloquée: la correction revient d'elle-même vers l'ancienne; le modèle à deux processus le reproduit, le modèle à un processus non. & \cite{Smith2006} & mesuré \\
15 & Chiffres de la fig.~\ref{fig:tworate}: $0.84$ (lent $0.63$) en $200$ mouvements; $6$ inverses; rapide $-0.53$, lent $0.46$; retour jusqu'à $0.37$ en $40$ mouvements & Calcul avec \texttt{models/\allowbreak motor\_\allowbreak learning.py}, $A_f=0.92$, $B_f=0.1$, $A_s=0.996$, $B_s=0.02$: $0.840$ ($0.634$ et $0.205$), $6$ mouvements, $-0.531$ et $0.457$, pic de $0.371$ au $40$\textsuperscript{e} mouvement. & \texttt{models/\allowbreak motor\_\allowbreak learning.py} & modèle \\
16 & Deux processus sont nécessaires parce que le corps change à des vitesses différentes & Théorie: l'estimation optimale sous des perturbations à plusieurs durées de vie donne plusieurs filtres de constantes de temps différentes et explique le retour spontané et le réapprentissage accéléré. & \cite{Kording2007} & hypothèse \\
\bottomrule
\end{longtable}}

\section{Chapitre~\ref{sec:chemoutput}. La sortie chimique}

Les deux fils vers le cœur et leurs vitesses différentes, la contre-réaction dans les cascades hormonales, le code par fréquence des bouffées, le générateur circadien et son recalage par la lumière sont mesurés directement; la limite $K<8$ est un théorème de l'automatique démontré dans le chapitre; les pulsations engendrées par la rétroaction de la cascade elle-même sont une hypothèse soutenue par une expérience solide.

{\footnotesize
\setlength{\tabcolsep}{3pt}\renewcommand{\arraystretch}{1.15}
\begin{longtable}{>{\raggedright}p{0.5cm}>{\raggedright}p{4.0cm}>{\raggedright}p{5.6cm}>{\raggedright}p{2.3cm}>{\raggedright\arraybackslash}p{1.7cm}}
\toprule
N\textsuperscript{o} & Affirmation & Expérience et ce qui est mesuré & Source & Statut \\
\midrule\endfirsthead
\toprule
N\textsuperscript{o} & Affirmation & Expérience et ce qui est mesuré & Source & Statut \\
\midrule\endhead
1 & Le stimulateur du cœur bat de lui-même environ $100$ fois par minute; au repos, le frein est enfoncé, et la fréquence est en général de l'ordre de $60$--$70$ & Humains, blocage complet des deux fils vers le cœur: la fréquence propre dépasse la fréquence de repos et décroît avec l'âge; chez l'adulte, elle est de l'ordre de $100$ battements par minute. Au repos, c'est donc le frein qui domine. La fréquence de repos de $60$--$70$ est une valeur typique, non citée séparément. & \cite{JoseCollison1970ihr} & mesuré \\
2 & L'accélérateur \nt{NORA} et le frein \nt{ACET} sont des filtres passe-bas, le frein est plus rapide~\eqref{eq:gasbrake} & Chien, stimulation modulée en fréquence du nerf vague ou du nerf sympathique cardiaque, et fonction de transfert vers la fréquence auriculaire: les deux voies sont des filtres passe-bas; la voie sympathique a une fréquence de coupure bien plus basse et un retard pur d'environ $1.7\,\text{s}$ en plus. Les caractéristiques dépendent du tonus moyen, si bien que la formule linéaire~\eqref{eq:gasbrake} est une approximation. & \cite{Berger1989} & mesuré \\
3 & Le frein est rapide parce qu'\nt{ACET} ouvre le canal potassique sans second messager, alors que \nt{NORA} passe par une chaîne de réactions & Fragments de membrane de cellules auriculaires: le canal potassique ouvert par l'acétylcholine est ouvert par les sous-unités $\beta\gamma$ de la protéine G couplée au récepteur, sans second messager intracellulaire. La voie de \nt{NORA} par l'AMPc n'est pas citée ici. & \cite{Logothetis1987gbg} & mesuré \\
4 & Le sang fait le tour du corps en une minute environ; \nt{ADRE} circulant atteint tous les organes munis du récepteur & Estimation générale d'ordre de grandeur; formulée ainsi dans le chapitre, sans source citée. & --- & estimation \\
5 & La cascade hormonale est bouclée par une contre-réaction: l'hormone finale freine les premiers étages & Revue d'expériences: les hormones corticosurrénaliennes répriment la libération d'ACTH par l'hypophyse et de l'hormone de libération par l'hypothalamus dans trois gammes de temps: rapide, intermédiaire et lente. & \cite{KellerWood1984fb} & mesuré \\
6 & Trois étages identiques sont stables pour $K<8$~\eqref{eq:cascadestable}; à la limite, la période vaut $2\pi\tau/\sqrt3\approx3.6\tau$ & Démontré dans le chapitre: à la pulsation $\sqrt3/\tau$, trois étages donnent un déphasage de $180^\circ$ et une atténuation d'un facteur $8$. & démonstration dans le chapitre & théorie \\
7 & Erreur sur le niveau $1/(1+K)$: un tel régulateur ne tient donc pas mieux qu'environ $10\,\%$ (proposition~\ref{prop:cascade}) & Conséquence de la ligne~6 pour une rétroaction proportionnelle. L'axe réel a une rétroaction sur plusieurs étages et dans plusieurs gammes de temps (ligne~5): la limite se rapporte donc au modèle~\eqref{eq:cascade} et n'est pas mesurée. & démonstration dans le chapitre & théorie \\
8 & Pour $K=4$ et $7$, le niveau se stabilise; pour $K=12$, pulsations régulières de période $3.7$--$3.9\,\tau$ (fig.~\ref{fig:cascade}) & Calcul avec \texttt{models/\allowbreak chem\_\allowbreak models.py}: pour $K=12$, périodes successives $3.9$, $3.7$, $3.8$, $3.7\,\tau$; pour $K=4$, amplitude en fin de calcul $0.003$. & \texttt{models/\allowbreak chem\_\allowbreak models.py} & modèle \\
9 & L'hormone du stress sort par pulsations environ une fois par heure; les pulsations naissent de la rétroaction elle-même, sans générateur séparé & Rats, perfusion constante de l'hormone de libération: l'ACTH et la corticostérone oscillent à une fréquence physiologique, proche d'une heure; une entrée rythmique venant de l'hypothalamus n'est pas nécessaire aux pulsations. Un modèle hypophyse--surrénale à rétroaction retardée produit de telles oscillations. & \cite{Walker2012glc, Walker2010} & hypothèse \\
10 & Le destinataire lit la fréquence des bouffées, non le niveau & Singes privés de leur propre libération de l'hormone de libération des gonadotrophines: une perfusion constante ne rétablit pas la sécrétion des hormones hypophysaires, des bouffées toutes les heures la rétablissent; le retour à la perfusion constante éteint de nouveau la réponse. La fatigue du récepteur vue comme filtre passe-haut est une interprétation. & \cite{Belchetz1978} & mesuré \\
11 & Des fréquences de bouffées différentes agissent différemment sur les diverses cellules d'une même glande & Mêmes singes: porter les bouffées à $2$--$5$ par heure fait baisser les deux hormones hypophysaires; les espacer à une toutes les $3$~heures fait baisser l'une (LH) et monter l'autre (FSH), si bien que leur rapport est fixé par la fréquence. & \cite{Wildt1981gnrh} & mesuré \\
12 & Le rythme circadien naît d'une contre-réaction intracellulaire retardée de plusieurs heures & Revue: les protéines des gènes d'horloge répriment avec retard la transcription de leurs propres gènes; la boucle transcription--traduction donne une période d'environ un jour. & \cite{TakahashiJS2017} & mesuré \\
13 & Le générateur principal est un groupe de quelques dizaines de milliers de neurones qui synchronise le reste du corps & Revue: le noyau suprachiasmatique est le générateur circadien principal; ses neurones isolés en culture sont des générateurs autonomes, et le réseau les synchronise; chez la souris, environ $10^4$ neurones de chaque côté. & \cite{Welsh2010scn} & mesuré \\
14 & La période propre chez l'humain est de $24.18$~heures & Humains en lumière faible, avec un horaire découplé du jour: période des rythmes de la mélatonine, de la température et du cortisol en moyenne de $24.18$~heures chez les jeunes comme chez les âgés, avec une faible dispersion. & \cite{Czeisler1999} & mesuré \\
15 & La lumière recale le générateur comme une boucle à verrouillage de phase~\eqref{eq:pll}; il faut un décalage d'environ $11$~min par jour & Humains, une seule exposition à une lumière vive de $6.7$~heures à différentes heures: courbe de réponse de phase de type~1, d'amplitude $5$~heures; retard avant le minimum de la température corporelle, avance après. L'équation~\eqref{eq:pll} est le modèle standard; $11$~min $=0.18$~heure est de l'arithmétique. & \cite{Khalsa2003prc} & mesuré \\
16 & Sans lumière, pas de verrouillage: le rythme dérive vers sa période propre & Personnes totalement aveugles, sans perception de la lumière, vivant en société normale: chez environ la moitié d'entre elles, les rythmes de la mélatonine et du cortisol suivent leur propre période, différente de celle du jour. & \cite{Sack1992blind} & mesuré \\
\bottomrule
\end{longtable}}

\section{Chapitre~\ref{sec:debugging}. Déboguer la machine}

Ce qu'entend une électrode, la faible dimension de l'activité, le fonctionnement des interfaces à grilles rigides, l'apprentissage conjoint du cerveau et du décodeur, et les points visuels produits par stimulation sont mesurés dans des expériences publiées avec comité de lecture; les estimations des flux de données sont l'arithmétique du chapitre; sur l'appareil de Neuralink implanté chez l'humain, pour autant que nous ayons pu le vérifier, les données ont surtout été publiées par la société elle-même.

{\footnotesize
\setlength{\tabcolsep}{3pt}\renewcommand{\arraystretch}{1.15}
\begin{longtable}{>{\raggedright}p{0.5cm}>{\raggedright}p{4.0cm}>{\raggedright}p{5.6cm}>{\raggedright}p{2.3cm}>{\raggedright\arraybackslash}p{1.7cm}}
\toprule
N\textsuperscript{o} & Affirmation & Expérience et ce qui est mesuré & Source & Statut \\
\midrule\endfirsthead
\toprule
N\textsuperscript{o} & Affirmation & Expérience et ce qui est mesuré & Source & Statut \\
\midrule\endhead
1 & L'électrode entend les impulsions des neurones situés dans un rayon d'une centaine de micromètres, en général d'un seul à quelques-uns; on les distingue par leur forme & Rat, région CA1, enregistrement intracellulaire et extracellulaire simultané: la forme de l'impulsion extracellulaire reproduit la largeur et l'amplitude de l'impulsion intracellulaire. Estimation: une tétrode pourrait distinguer $60$--$100$ neurones, alors qu'en pratique on en isole environ six. & \cite{Henze2000extra} & mesuré \\
2 & Le cerveau compte $8.6\cdot10^{10}$ neurones; la sonde en entend environ un cent-millionième & Comptage des noyaux dans une suspension de tissu dissous: $86\cdot10^9$ neurones dans le cerveau humain. La fraction est l'arithmétique du chapitre. & \cite{Azevedo2009} & mesuré \\
3 & L'activité de centaines de neurones du cortex moteur tient dans une dizaine ou deux de variables communes & Revue d'enregistrements dans le cortex moteur du singe: l'activité de la population est décrite par quelques variables latentes communes à de nombreux neurones. & \cite{Gallego2017} & mesuré \\
4 & Le bruit des voisins est corrélé, et une centaine de neurones porte presque autant qu'un millier (proposition~\ref{prop:pooling}) & Cortex visuel du singe: coefficient de corrélation du bruit entre neurones voisins d'environ $0.1$, et le gain de la moyenne sature vers une centaine de neurones. Théorie des corrélations qui limitent l'information. & \cite{Zohary1994, MorenoBote2014} & mesuré \\
5 & Flux brut de $2\cdot10^8$~bit/s contre $8\cdot10^4$~bit/s en impulsions~\eqref{eq:probedata} & Arithmétique du chapitre à partir de la capacité du flux d'impulsions~\eqref{eq:spikeentropy}. Les paramètres de numérisation sont réels: la puce de Neuralink échantillonne à $19.3$~kHz sur $10$~bits par canal. & démonstration dans le chapitre; \cite{Rieke1997, Musk2019} & théorie \\
6 & Premier système Neuralink: les puces amplifient et numérisent, le flux brut part par câble, un programme externe détecte les impulsions; détection sur puce, boîtier fermé, radio et alimentation sans fil: pour l'appareil destiné à l'humain, selon la société & Article de la société: tout le flux brut passe par un câble USB-C, et un programme détecte les impulsions en temps réel hors de la puce; compression embarquée, alimentation et transmission sans fil sont présentées comme un projet pour les appareils cliniques. Pour l'appareil implanté chez l'humain, seulement des communications de la société. Que la détection des impulsions sur puce soit nécessaire est une conclusion du chapitre tirée de~\eqref{eq:probedata}. & \cite{Musk2019}; rapports de Neuralink, pas d'article à comité de lecture & mesuré (article de la société); chez l'humain: rapport de la société \\
7 & Jusqu'à $3072$ électrodes sur $96$ fils de $32$; un robot insère les fils en évitant les vaisseaux & Article de la société: $3072$ électrodes sur $96$ fils de $32$; le robot insère $6$ fils ($192$ électrodes) par minute; chez le rat porteur d'une grille implantée à long terme, jusqu'à $70\,\%$ des électrodes entendent des impulsions. & \cite{Musk2019} & mesuré (article de la société) \\
8 & Chez l'humain, environ mille canaux sur $64$ fils; une partie des fils s'est déplacée; curseur de l'ordre de $10$~bit/s & Communications de la société seulement. Une recherche dans PubMed n'a trouvé aucun article à comité de lecture présentant des résultats chez l'humain, ce qui concorde avec la réserve faite dans le texte du chapitre. & rapports de Neuralink; pas d'article à comité de lecture & mesuré (rapport de la société) \\
9 & Une interface sur une grille d'une centaine d'électrodes pilote un curseur & Personne paralysée, $96$ électrodes dans le cortex moteur primaire: l'intention de bouger le bras modifie les impulsions trois ans après la lésion; le décodeur fournit un \og curseur neuronal\fg{} (courrier, télévision), la commande d'une prothèse de main et d'un bras robotisé. & \cite{Hochberg2006} & mesuré \\
10 & Écriture \og par la main mentale\fg{}: environ $90$ caractères par minute, moins de $8$~bit/s & Personne à la main paralysée, tentatives d'écriture manuscrite, décodeur récurrent: $90$ caractères par minute avec une précision de $94.1\,\%$ en temps réel, plus de $99\,\%$ après correction automatique. L'estimation en bits est l'arithmétique du chapitre. & \cite{Willett2021} & mesuré \\
11 & Les tentatives de parole sont converties en texte à environ $60$ mots par minute & Personne ayant perdu une parole intelligible, grilles dans le cortex: $62$ mots par minute; $9.1\,\%$ d'erreurs sur les mots avec un vocabulaire de $50$ mots, $23.8\,\%$ avec un vocabulaire de $125\,000$. & \cite{Willett2023} & mesuré \\
12 & L'interface n'extrait qu'une petite part de la capacité: quelques dizaines de bits par seconde pour un neurone, quelques milliers pour une centaine (proposition~\ref{prop:probe}) & La borne supérieure~\eqref{eq:spikeentropy} est théorique; la comparaison avec les lignes~8 et~10 est une conclusion du chapitre. Que le goulot d'étranglement soit la faible dimension et la méconnaissance du code, et non le fil, n'a pas été vérifié par une expérience directe. & \cite{Rieke1997} & théorie \\
13 & Le cerveau et le décodeur apprennent l'un de l'autre & Singes pilotant un curseur; le décodeur s'adapte en boucle fermée pendant que les neurones réapprennent: la précision augmente, l'habileté se conserve et souffre moins des mouvements du bras propre. Le conseil de séparer les vitesses d'apprentissage est une règle d'ingénieur, sans expérience propre dans le chapitre. & \cite{Orsborn2014coad} & mesuré \\
14 & Le courant envoyé par une électrode excite les neurones et les fils environnants sans distinction de type & Cortex de souris et de chat, imagerie calcique à deux photons: même un faible courant excite des neurones épars, dispersés jusqu'à des millimètres, par excitation directe des axones dans un volume de quelques dizaines de micromètres de diamètre. L'excitation est donc non sélective, mais pas continue. & \cite{Histed2009stim} & mesuré \\
15 & La stimulation du cortex visuel allume des points; jusqu'à $16$ points à la fois permettent de reconnaître certaines lettres et les contours des objets & Personne aveugle, $96$ électrodes dans le cortex visuel pendant $6$ mois: seuil d'une électrode $66.8\pm36.5$~µA; la stimulation simultanée de groupes (jusqu'à $16$ électrodes, limite du stimulateur) abaisse le seuil et produit des motifs distincts; certaines lettres et les contours des objets sont reconnus. & \cite{Fernandez2021} & mesuré \\
16 & Le second axe de Neuralink est un appareil qui envoie des signaux dans le cortex visuel & Seulement des communications de la société sur son développement; aucune donnée à comité de lecture sur son fonctionnement n'a été trouvée. & rapports de Neuralink & hypothèse \\
17 & On peut mesurer les concentrations des neurotransmetteurs de diffusion par un capteur chimique dans le tissu & Tranches de cerveau de rat, voltampérométrie cyclique rapide sur fibre de carbone: libération et recapture de la sérotonine mesurées; la substance sort de la synapse. & \cite{Bunin1998} & mesuré \\
\bottomrule
\end{longtable}}

\section{Chapitre~\ref{sec:eetypes}. Les neurones comme composants}

Les modes de fonctionnement des cellules individuelles sont mesurés directement, par un courant injecté par électrode et un enregistrement de la tension; les mathématiques de l'intégrateur et de la division en classes relèvent de la théorie classique, tandis que l'usage par le cerveau de la reconfiguration des modes pour calculer reste une hypothèse.

{\footnotesize
\setlength{\tabcolsep}{3pt}\renewcommand{\arraystretch}{1.15}
\begin{longtable}{>{\raggedright}p{0.5cm}>{\raggedright}p{4.0cm}>{\raggedright}p{5.6cm}>{\raggedright}p{2.3cm}>{\raggedright\arraybackslash}p{1.7cm}}
\toprule
N\textsuperscript{o} & Affirmation & Expérience et ce qui est mesuré & Source & Statut \\
\midrule\endfirsthead
\toprule
N\textsuperscript{o} & Affirmation & Expérience et ce qui est mesuré & Source & Statut \\
\midrule\endhead
1 & Résonateur: un courant lent qui s'oppose aux variations de tension fait du neurone un filtre passe-bande dont le maximum se situe entre quelques hertz et quelques dizaines de hertz (section~\ref{sec:resonator}) & En tranches, on injectait dans des neurones \nt{GLUT} un courant sinusoïdal de fréquence croissante et on mesurait l'impédance $|Z(f)|$: maximum à $2$--$5$~Hz à $33^\circ$C (environ $7$~Hz à $38^\circ$C); la résonance est créée par des courants lents potassique et cationique, et amplifiée par le courant sodique persistant. Une revue décrit la même méthode pour de nombreuses classes de cellules & \cite{HuVervaekeStorm2002,HutcheonYarom2000} & mesuré \\
2 & Le courant lent se comporte comme une inductance; avec la capacité de la membrane, il forme un circuit oscillant & La réponse sous le seuil d'un axone à un courant a été mesurée et décrite par des équations de conductance linéarisées; le schéma équivalent contient une inductance \og phénoménologique\fg{} dont la valeur dépend de la tension & \cite{Mauro1970} & théorie \\
3 & Constante de temps d'un groupe à rétroaction $\tau_{\text{eff}}=\tau/(1-w)$~\eqref{eq:perfectint}; pour $\tau=0.1$~s et $w=0.995$, $20$~s & Déduite dans le chapitre de l'équation linéaire du groupe; proposée comme mécanisme de maintien de la position de l'œil dans un travail théorique & démonstration dans le chapitre; \cite{Seung1996} & théorie \\
4 & L'intégrateur de position de l'œil maintient son entrée grâce au réseau, et non grâce à une propriété d'une cellule isolée & Enregistrement intracellulaire chez le poisson, dans le noyau qui maintient la position de l'œil: après une rotation rapide, la fréquence du neurone change par palier et se maintient; un bref courant injecté dans une seule cellule ne provoque qu'un décalage passager, tandis que les paliers de tension persistent même sous le seuil: ce sont les entrées synaptiques qui les entretiennent & \cite{Aksay2001} & mesuré \\
5 & L'intégrateur sans fuite a besoin d'un réglage permanent par apprentissage; déréglé, il oublie ou part en saturation & On entraînait des poissons avec une rétroaction visuelle proportionnelle à $\pm$ la position de l'œil: le maintien devenait instable ou fuyant, et la constante de temps des neurones chutait de plus d'un ordre de grandeur, jusqu'à $<1$~s; la rétroaction visuelle normale rétablissait peu à peu la stabilité & \cite{Major2004} & mesuré \\
6 & Neurone bistable: une brève entrée enclenche un plateau durable, l'inhibition l'arrête; la propriété est activée par la sérotonine (section~\ref{sec:bistable}) & Enregistrement intracellulaire chez le chat de neurones \nt{ACET} qui commandent les muscles: une brève excitation synaptique fait passer le neurone en activité durable, une brève inhibition le remet à zéro; chez le chat spinal, l'état apparaît après administration d'un précurseur de la sérotonine & \cite{Hounsgaard1988} & mesuré \\
7 & Deux classes: les intégrateurs (fréquence croissant à partir de zéro) et les résonateurs (fréquence finie d'emblée au seuil) & Classes déduites de la théorie des bifurcations. Expérience: en tranches de cortex, les neurones \nt{GLUT} se mettent à décharger à une fréquence aussi basse qu'on veut (type~1), les neurones \nt{GABA} rapides démarrent d'un bond à fréquence finie (type~2) & \cite{Izhikevich2007,Tateno2004} & mesuré \\
8 & Un même neurone est intégrateur ou détecteur de coïncidences: l'inhibition raccourcit la fenêtre de sommation & En tranches d'hippocampe, la fenêtre dans laquelle les entrées excitatrices s'additionnent jusqu'au seuil est inférieure à $2$~ms; elle est raccourcie par l'inhibition directe qui suit immédiatement l'excitation; dans les dendrites, où l'inhibition est plus faible, la fenêtre est plus large & \cite{Pouille2001} & mesuré \\
9 & Ligne à retard faite d'axones (tableau~\ref{tab:eetypes}) & Chez la chouette effraie, on a mesuré les retards dans les axones qui vont aux détecteurs de coïncidences: des longueurs d'axone différentes donnent des retards différents, et les détecteurs sont accordés sur la différence des temps d'arrivée du son & \cite{Carr1990} & mesuré \\
10 & Stimulateur pendant la veille, cellule silencieuse pendant le sommeil & Les neurones \nt{SERO} déchargent de façon presque régulière le jour, ralentissent en sommeil lent et se taisent presque en sommeil paradoxal (détails au chapitre~\ref{sec:sleep}) & \cite{JacobsAzmitia1992,McGintyHarper1976} & mesuré \\
11 & Le composant est fixé par les canaux ioniques et par les canaux de diffusion qui les ajustent (proposition~\ref{prop:eetypes}) & Montré sur de petits réseaux: les substances régulatrices modifient les propriétés propres des neurones et la force des synapses, si bien qu'un même circuit produit des rythmes différents & \cite{Marder2012} & mesuré \\
12 & Le cerveau exploite la reconfiguration des modes pour calculer (proposition~\ref{prop:eetypes}) & Il n'existe aucune expérience directe où la reconfiguration du mode d'une cellule serait nécessaire pour résoudre une tâche; seulement des expériences où les régulateurs modifient la sortie du circuit & \cite{Marder2012} & hypothèse \\
\bottomrule
\end{longtable}}

\section{Chapitre~\ref{sec:dendrites}. Le nombre de dendrites}

La structure des classes et le nombre d'entrées sont mesurés par l'anatomie et par des enregistrements électriques; l'estimation~\eqref{eq:dendriteload} relève de la simple physique du condensateur, tandis que l'explication computationnelle du nombre d'entrées s'appuie sur la théorie et sur des modèles.

{\footnotesize
\setlength{\tabcolsep}{3pt}\renewcommand{\arraystretch}{1.15}
\begin{longtable}{>{\raggedright}p{0.5cm}>{\raggedright}p{4.0cm}>{\raggedright}p{5.6cm}>{\raggedright}p{2.3cm}>{\raggedright\arraybackslash}p{1.7cm}}
\toprule
N\textsuperscript{o} & Affirmation & Expérience et ce qui est mesuré & Source & Statut \\
\midrule\endfirsthead
\toprule
N\textsuperscript{o} & Affirmation & Expérience et ce qui est mesuré & Source & Statut \\
\midrule\endhead
1 & Capacité spécifique de la membrane $c_m\approx1$~µF/cm$^2$ & On prélevait la membrane du corps d'un neurone dans une pipette, on appliquait un échelon de tension et, d'après la décroissance du courant capacitif, on trouvait la capacité totale, puis on la divisait par la surface: $0.9$~µF/cm$^2$ pour trois classes de neurones & \cite{Gentet2000} & mesuré \\
2 & Temps de charge $t\approx c_mA\Delta V/I$~\eqref{eq:dendriteload}: un grand neurone a besoin de dizaines d'entrées simultanées, une seule grosse synapse suffit à un petit & Estimation par la loi de charge d'un condensateur; les nombres ($20$ et $300$~pF, $0.1$ et quelques nA) sont des ordres de grandeur & démonstration dans le chapitre & théorie \\
3 & Une seule synapse énorme délivre un courant de quelques nanoampères et amène aussitôt un petit neurone au seuil & Enregistrement simultané de la terminaison et du destinataire en tranche: courant d'une synapse de $2$--$13$~nA, chaque impulsion d'entrée provoque une réponse suprasliminaire, avec un délai d'environ $0.4$~ms à $36^\circ$C; le destinataire émet \nt{GLYC} & \cite{Borst1995,BorstSoria2012} & mesuré \\
4 & Zéro dendrite: dans les neurones sensoriels du toucher et de la douleur, l'impulsion va du capteur à la moelle épinière sans passer par le corps cellulaire & La structure (un axone qui se divise en deux près du corps) est établie par l'anatomie. Que la conduction n'exige pas l'excitabilité du corps a été montré sur un modèle validé d'un tel neurone: sans corps excitable, l'impulsion franchit la bifurcation avec la même fiabilité & \cite{AmirDevor2003} & théorie \\
5 & Une dendrite: le neurone olfactif porte un seul type de récepteur et transmet une seule ligne d'entrée & Revue d'expériences sur les gènes des récepteurs: chaque neurone olfactif exprime un seul gène de récepteur parmi une grande famille & \cite{Mombaerts2004} & mesuré \\
6 & Deux dendrites: dans les détecteurs de direction du son, les entrées des oreilles gauche et droite se posent sur des dendrites différentes & Anatomie et enregistrements chez les mammifères, réunis dans une revue: les entrées des deux oreilles sont réparties sur deux dendrites opposées & \cite{Grothe2010} & mesuré \\
7 & Répartir les entrées sur deux dendrites rend le \og ET\fg{} plus strict & Montré sur un modèle: la saturation~\eqref{eq:shunt} sur une seule dendrite empêche les entrées d'un même côté d'atteindre le seuil, et la détection des coïncidences s'améliore. Il n'existe pas d'expérience directe où l'on aurait modifié la disposition des entrées & \cite{AgmonSnir1998} & théorie \\
8 & Environ $7\cdot10^{10}$ petits neurones \nt{GLUT} dans le circuit d'apprentissage moteur, avec $\sim4$ entrées chacun & Comptage des cellules par fractionnement isotrope: environ $69$~milliards de neurones dans cette partie du cerveau humain. Enregistrement d'une telle cellule chez le rat vivant: un contact déclenche des bouffées de courants discrets venant de quelques entrées, et la cellule décharge quand ils s'additionnent & \cite{Azevedo2009,Chadderton2004} & mesuré \\
9 & Un élément à seuil à $n$ entrées sépare au plus $P_{\max}\approx2n$ motifs aléatoires~\eqref{eq:cover}; pour le neurone de sortie du circuit d'apprentissage moteur, la borne est $\sim2\cdot10^5$, et l'estimation réaliste $\sim4\cdot10^4$ associations & La borne est un résultat classique de géométrie combinatoire. L'estimation réaliste vient d'une théorie de la capacité avec des poids uniquement positifs et une marge contre le bruit, appliquée au nombre mesuré d'entrées de ce neurone & \cite{Cover1965,Brunel2004} & théorie \\
10 & Les mélangeurs à peu d'entrées servent à déployer l'entrée; \og quatre\fg{} est proche de l'optimum & Théorie: la dimension de la représentation que fournit une couche de mélangeurs est maximale à peu près pour le nombre d'entrées observé en anatomie; l'hypothèse initiale de l'expansion figure dans des travaux classiques & \cite{Marr1969,Albus1971,LitwinKumar2017} & hypothèse \\
11 & Grands arbres: $10^4$--$10^5$ entrées & Comptage des synapses sur des images de microscopie électronique: $\approx1.7\cdot10^5$ synapses sur un seul neurone \nt{GABA} de sortie du circuit d'apprentissage moteur chez le rat; environ $30\,000$ entrées excitatrices et $1\,700$ inhibitrices sur un neurone \nt{GLUT} de l'hippocampe & \cite{NapperHarvey1988,Megias2001} & mesuré \\
12 & Les branches d'un grand arbre sont des sous-circuits distincts dotés de leurs propres seuils; le neurone est un petit réseau multicouche & Stimulation ponctuelle de deux sites sur de fines dendrites: les entrées d'une même branche s'additionnent selon une sigmoïde, celles de branches différentes linéairement. Dans les dendrites de neurones du cortex humain, on a trouvé des impulsions calciques graduées qui permettent à une seule cellule de résoudre un problème non linéairement séparable. Modèle: pour reproduire un seul neurone, il faut un réseau profond & \cite{Polsky2004,Gidon2020,Beniaguev2021} & mesuré \\
13 & Dendrites-sorties: chaque branche d'un neurone \nt{GABA} de la rétine est un détecteur de direction distinct & Imagerie calcique à deux photons dans des branches isolées: le signal d'une branche est sélectif au mouvement allant du corps cellulaire vers la pointe, alors que la tension du corps ne l'est pas & \cite{Euler2002} & mesuré \\
14 & Le nombre d'entrées suit la tâche (proposition~\ref{prop:dendrites}) & La structure des classes est mesurée (lignes 3--13); que le nombre d'entrées soit choisi pour la vitesse ou pour la classification n'a été montré que par la théorie et des modèles, pour certaines classes & \cite{LitwinKumar2017,AgmonSnir1998} & hypothèse \\
\bottomrule
\end{longtable}}

\section{Chapitre~\ref{sec:sleep}. Le sommeil}

Les modes de sommeil, les répétitions et la diminution des synapses sont mesurés par des enregistrements de neurones, par microdialyse et par microscopie électronique; l'horaire du sommeil et la bascule lente sont décrits par les modèles du livre, tandis que la fonction du sommeil relève surtout d'hypothèses.

{\footnotesize
\setlength{\tabcolsep}{3pt}\renewcommand{\arraystretch}{1.15}
\begin{longtable}{>{\raggedright}p{0.5cm}>{\raggedright}p{4.0cm}>{\raggedright}p{5.6cm}>{\raggedright}p{2.3cm}>{\raggedright\arraybackslash}p{1.7cm}}
\toprule
N\textsuperscript{o} & Affirmation & Expérience et ce qui est mesuré & Source & Statut \\
\midrule\endfirsthead
\toprule
N\textsuperscript{o} & Affirmation & Expérience et ce qui est mesuré & Source & Statut \\
\midrule\endhead
1 & Pendant le sommeil, les neurones du cortex ne se taisent pas; c'est le mode de fonctionnement qui change & Longs enregistrements de neurones du cortex chez le rat: en sommeil lent, la fréquence reste non nulle, et des périodes d'activité alternent avec des périodes de silence; après une longue veille, la fréquence est plus élevée dans tous les états, après le sommeil elle est plus basse & \cite{Vyazovskiy2009} & mesuré \\
2 & \nt{ACET} est élevé le jour et en sommeil paradoxal, bas en sommeil lent (tableau~\ref{tab:sleepmodes}) & Microdialyse dans le cortex chez des chats se déplaçant librement: pendant la veille, la libération d'acétylcholine est supérieure de $100$--$175\%$ à celle du sommeil lent, et en sommeil paradoxal elle est à peu près au niveau de la veille calme & \cite{Marrosu1995} & mesuré \\
3 & \nt{NORA} et \nt{SERO} s'apaisent en sommeil lent et se taisent presque en sommeil paradoxal & Enregistrement de neurones \nt{NORA} isolés chez le rat et de neurones \nt{SERO} chez le chat selon les stades du sommeil: la fréquence est maximale pendant la veille, plus basse en sommeil lent et presque nulle en sommeil paradoxal & \cite{AstonJonesBloom1981,McGintyHarper1976} & mesuré \\
4 & En sommeil paradoxal, les muscles sont coupés par une inhibition via \nt{GABA} et \nt{GLYC} & Chez le rat, on mesurait l'activité des neurones \nt{ACET} du muscle masticateur et on leur appliquait directement des bloqueurs: la paralysie n'est levée que si l'on bloque à la fois les deux types de récepteurs \nt{GABA} et les récepteurs \nt{GLYC} & \cite{BrooksPeever2012} & mesuré \\
5 & La pression de sommeil $S$ est un condensateur à deux seuils~\eqref{eq:sleeppressure}, $\tau_r=18.2$~h, $\tau_d=4.2$~h & Modèle à deux processus; les constantes de temps sont tirées de la façon dont la puissance des ondes lentes de l'EEG croît avec la veille et décroît pendant le sommeil, et la forme des seuils de l'heure du réveil spontané & \cite{Borbely1982,Daan1984} & théorie \\
6 & La machine trouve d'elle-même $16.1$~h de veille et $8.0$~h de sommeil (fig.~\ref{fig:twoprocess}) & Calcul selon~\eqref{eq:sleeppressure} avec des seuils oscillants, script \texttt{models/sleep\_models.py}: $16.1$~h de veille et $7.9$--$8.0$~h de sommeil & --- & modèle \\
7 & Après $40$~h sans sommeil, $S=0.90$, mais le sommeil ne dure que $8.8$~h: la dette se rembourse par la profondeur, et non par la durée & Modèle: \texttt{models/sleep\_models.py} donne $S=0.90$ et $8.8$~h. Expérience: après $40.5$~h sans sommeil, chez huit personnes, la puissance des ondes lentes de l'EEG est significativement plus élevée que d'habitude lors de la première nuit de récupération & \cite{Borbely1981} & modèle \\
8 & Physiquement, $S$ est l'adénosine & Microdialyse chez le chat: l'adénosine extracellulaire, dans l'une des régions qui commandent la veille, augmente pendant la veille, augmente encore en cas de privation de sommeil et baisse pendant le sommeil. Que $S$ soit uniquement l'adénosine n'est pas démontré & \cite{PorkkaHeiskanen1997,Dunwiddie2001} & hypothèse \\
9 & En sommeil lent, le cortex bascule: quelques centaines de ms d'activité, puis le silence, moins d'une fois par seconde ($<1$~Hz, le plus souvent $0.3$--$0.4$~Hz sous anesthésie) & Enregistrement intracellulaire chez le chat anesthésié: dépolarisation de $0.8$--$1.5$~s et longue hyperpolarisation, rythme $<1$~Hz, le plus souvent $0.3$--$0.4$~Hz; la même alternance d'activité et de silence s'observe dans le sommeil naturel (ligne~1) & \cite{Steriade1993,Vyazovskiy2009} & mesuré \\
10 & La bascule est un groupe à rétroaction et fatigue~\eqref{eq:updown}; pour $b=5$, période de $1.1$~s (valeur du modèle) et activité environ $35\%$ du temps; pour $b=1$, fonctionnement régulier (fig.~\ref{fig:updown}) & Calcul, script \texttt{models/sleep\_models.py}: période $1.11$~s, fraction du temps d'activité $0.35$ (moyenne sur un nombre entier de périodes); pour $b=1$, fraction d'activité $1.0$. La période du modèle est plus courte que la période mesurée (ligne~9) & --- & modèle \\
11 & \nt{ACET} arrête la bascule et fait passer le cortex à un fonctionnement régulier & La stimulation d'un noyau de neurones \nt{ACET} chez le chat bloquait le rythme lent dans $79\%$ des neurones du cortex et les faisait passer à un fonctionnement régulier, surtout par disparition des longues hyperpolarisations. L'acétylcholine supprime les courants potassiques lents qui créent la fatigue & \cite{SteriadeAmzica1993,McCormick1992} & mesuré \\
12 & Les bouffées de rythme à $7$--$15$~Hz naissent d'une boucle \nt{GLUT}--\nt{GABA} entre le cortex et un nœud relais & Dans des tranches du nœud relais visuel du furet, les bouffées de rythme sont créées par les bouffées de réponse des cellules relais à l'inhibition venant d'un noyau \nt{GABA} voisin, qu'elles excitent elles-mêmes & \cite{vonKrosigk1993,SteriadeMcCormickSejnowski1993} & mesuré \\
13 & Les répétitions de la journée se font en accéléré: environ $20$ fois plus vite dans l'hippocampe, $6$--$7$ fois dans le cortex & En sommeil lent, les séquences répétées de neurones de l'hippocampe sont comprimées dans le temps. Après une course sur une piste, les séquences de cellules de lieu se répétaient dans le même ordre par brèves bouffées de $\sim100$~ms, comprimées environ $20$ fois; dans le cortex, compression de $6$--$7$ fois & \cite{Nadasdy1999,LeeWilson2002,Euston2007} & mesuré \\
14 & Renforcer la coordination des répétitions avec le cortex améliore la mémoire (lien causal) & Chez le rat, après l'apprentissage, une stimulation calée sur les répétitions renforçait leur couplage avec les ondes lentes et les bouffées de rythme du cortex: le lendemain, la mémoire était élevée, alors que chez les témoins elle restait au niveau du hasard & \cite{Maingret2016} & mesuré \\
15 & La fonction des répétitions est l'apprentissage entrelacé de la mémoire lente & Théorie des deux systèmes d'apprentissage et calculs sur des réseaux artificiels: l'apprentissage séquentiel abîme l'ancien, l'apprentissage entrelacé non. Il n'existe pas d'expérience directe sur le cerveau montrant que les répétitions servent précisément à cela & \cite{McClelland1995} & hypothèse \\
16 & Après le sommeil, les synapses diminuent proportionnellement à leur taille; les grosses changent à peine & Microscopie électronique tridimensionnelle de $6920$ synapses du cortex de souris: surface de contact inférieure de $\sim18\%$ après le sommeil par rapport à la veille, diminution proportionnelle à la taille, les grosses synapses stables ne changent pas & \cite{deVivo2017} & mesuré \\
17 & La tâche principale du sommeil est la renormalisation des poids & Hypothèse de l'homéostasie synaptique; les lignes 1 et 16 s'accordent avec elle, mais ne prouvent pas qu'il s'agit de la fonction principale & \cite{TononiCirelli2014} & hypothèse \\
18 & Le sommeil paradoxal est un essai sur banc du modèle du monde & Le mode (tableau~\ref{tab:sleepmodes}) est mesuré; que le réseau fasse tourner un modèle du monde est une supposition théorique, sans expérience & \cite{Hobson2009} & hypothèse \\
19 & Lavage du cerveau pendant le sommeil: question ouverte & Une expérience: pendant le sommeil, l'espace intercellulaire est $60\%$ plus large et le nettoyage plus rapide. Une autre, avec une autre méthode de mesure: pendant le sommeil et sous anesthésie, le nettoyage est nettement plus lent. Les résultats se contredisent & \cite{Xie2013,Miao2024} & hypothèse \\
20 & Sommeil local: des zones du cortex d'un animal éveillé sombrent brièvement dans le silence & Chez le rat, après une longue veille, les neurones d'une zone du cortex se taisent brièvement, comme en sommeil lent, avec une onde lente dans l'EEG local; à ces moments-là, le rat se trompe plus souvent dans la tâche & \cite{Vyazovskiy2011} & mesuré \\
\bottomrule
\end{longtable}}

\section{Chapitre~\ref{sec:livingmachine}. Une machine faite de neurones vivants}

Le comportement des cultures sur matrices d'électrodes est mesuré dans des expériences directes d'enregistrement et de stimulation; le mécanisme des salves s'appuie sur le modèle d'épuisement synaptique et sur des expériences en microcultures, et les affirmations sur la dopamine en boîte de culture sur des expériences isolées, dont certaines ne sont, de l'aveu même des auteurs, qu'une démonstration.

{\footnotesize
\setlength{\tabcolsep}{3pt}\renewcommand{\arraystretch}{1.15}
\begin{longtable}{>{\raggedright}p{0.5cm}>{\raggedright}p{4.0cm}>{\raggedright}p{5.6cm}>{\raggedright}p{2.3cm}>{\raggedright\arraybackslash}p{1.7cm}}
\toprule
N\textsuperscript{o} & Affirmation & Expérience et ce qui est mesuré & Source & Statut \\
\midrule\endfirsthead
\toprule
N\textsuperscript{o} & Affirmation & Expérience et ce qui est mesuré & Source & Statut \\
\midrule\endhead
1 & Culture sur puce: surtout des \nt{GLUT}, une part plus faible de \nt{GABA} qui dépend de la préparation; environ $6\%$ dans les cultures d'hippocampe & Un réseau de milliers de neurones sur matrice d'électrodes est une expérience standard (ligne~3). La part des neurones \nt{GABA} dans les cultures d'hippocampe a été comptée par marquage du GABA: environ $6\%$ des neurones. Pour les cultures de cortex, nous ne donnons pas de chiffre vérifié & \cite{Benson1994} & mesuré \\
2 & Organoïdes: tissu nerveux tridimensionnel d'environ un millimètre, issu de cellules souches humaines & À partir de cellules souches, on a cultivé des organoïdes tridimensionnels comportant plusieurs régions cérébrales, dont un cortex avec des couches de progéniteurs et des neurones matures & \cite{Lancaster2013} & mesuré \\
3 & Sans entrée, la culture s'embrase en salves espacées d'une seconde à plusieurs minutes, selon la culture & $58$ cultures de cortex, d'une densité de $3\,000$ à $50\,000$ neurones, ont été enregistrées pendant cinq semaines ($963$ enregistrements d'une demi-heure): après deux semaines, les salves de tout le réseau dominent dans presque toutes les cultures; les intervalles entre salves vont de $1$ à $300$~s et dépendent fortement de la préparation & \cite{Wagenaar2006} & mesuré \\
4 & Une salve s'achève par épuisement du neurotransmetteur, avec un intervalle $T_{\text{salve}}\approx\tau_{\text{rec}}\ln\frac{1}{1-x^*}$~\eqref{eq:burstinterval}; $2$--$4$~s est l'estimation du modèle avec $\tau_{\text{rec}}=0.8$~s, la récupération mesurée est plus lente & La formule est établie dans le chapitre. Modèle: un réseau à synapses dépressives produit spontanément des salves de tout le réseau. Expérience sur des microcultures de $4$--$30$ neurones: la durée d'une salve dépend du temps écoulé depuis la précédente, et l'épuisement des vésicules de neurotransmetteur supprime les salves; conclusion des auteurs: la salve est limitée par la réserve de neurotransmetteur. La réserve s'y reconstitue en quelques dizaines de secondes, ce qui, avec la même formule, donne des intervalles de dizaines de secondes à plusieurs minutes & démonstration dans le chapitre; \cite{Tsodyks2000,CohenSegal2011,Tsodyks1997} & théorie \\
5 & \og Le professeur qui se tait\fg{}: la stimulation est coupée quand la réponse voulue apparaît, et la réponse se fixe & La culture était stimulée à $0.3$--$1$~Hz jusqu'à ce que la réponse choisie apparaisse $50\pm10$~ms après le stimulus, puis la stimulation était coupée pendant $5$~min; après quelques cycles, la réponse voulue était évoquée d'emblée; des courbes d'apprentissage ont été tracées & \cite{Shahaf2001} & mesuré \\
6 & La culture joue au tennis; hausse significative des séries en une session seulement avec rétroaction complète & Détails au chapitre~\ref{sec:dishbrain} et dans son annexe: avec du silence à la place du stimulus, la culture joue aussi mieux en fin de session que sans rétroaction, mais avec un effet moindre; l'apprentissage d'une session à l'autre est instable & \cite{Kagan2022} & mesuré \\
7 & Que la culture apprenne à prédire son entrée reste une hypothèse; les grands mots sur sa \og sensibilité\fg{} sont contestés & Le principe de l'énergie libre est une proposition théorique; aucune expérience directe ne le distingue d'explications plus simples. Dans un article de réponse, trente chercheurs ont souligné qu'un changement de comportement dans la boucle ne démontre ni intelligence ni sensations & \cite{Friston2010,Balci2023} & hypothèse \\
8 & La culture dirige un corps virtuel vers une cible avec un motif de stimuli bien choisi & Un programme choisissait lui-même les motifs de stimulation d'entraînement selon le résultat courant; en quelques dizaines de minutes, le réseau atteignait les états visés, et la plasticité durait $80$~min après $2$~h d'entraînement. Les mêmes stimuli, appliqués sans lien avec le comportement, restaient sans effet & \cite{Bakkum2008} & mesuré \\
9 & Réservoir: seule la lecture linéaire $y=W_{\text{sort}}\mathbf r$~\eqref{eq:reservoir} est entraînée & Théorie du calcul par réservoir: tout système non linéaire à mémoire évanescente, plus une sortie linéaire entraînée & \cite{Maass2002,JaegerHaas2004} & théorie \\
10 & Un organoïde utilisé comme réservoir distingue les locuteurs avec une précision d'environ $78\%$ (selon les auteurs); la lecture apprend par l'erreur, et les connexions de l'organoïde changent sans professeur & Des sons de parole de plusieurs locuteurs étaient injectés à l'organoïde sous forme de motifs de courant sur une matrice d'électrodes; la lecture entraînée distinguait le locuteur avec une précision d'environ $78\%$, selon les auteurs. Ceux-ci rapportent aussi que la connectivité fonctionnelle de l'organoïde changeait pendant l'entraînement, et parlent d'apprentissage non supervisé dans l'organoïde lui-même & \cite{Cai2023} & mesuré \\
11 & Un million de neurones dépense environ $10^{-4}$~W en impulsions~\eqref{eq:dishpower} & Estimation: le coût d'une impulsion, $10^{-10}$~J, tiré du budget énergétique du cortex~\eqref{eq:energy}, multiplié par le nombre d'impulsions; il n'existe pas de mesure directe de la puissance d'une culture & démonstration dans le chapitre; \cite{Attwell2001} & théorie \\
12 & Dopamine par éclair de lumière: $1$~mM de dopamine en cage, $240$ répétitions, le nombre d'impulsions évoquées augmentait & Sur une plateforme d'organoïdes à distance, la dopamine enfermée dans une cage photosensible ($1$~mM) était libérée par ultraviolet ($365$~nm, $0.8$~s) quand le stimulus provoquait plus d'impulsions; en $240$ étapes (environ $5$~h), le nombre d'impulsions a globalement augmenté. Les auteurs écrivent qu'une seule expérience ne permet pas d'établir le lien avec la dopamine; il n'y a pas de contrôles & \cite{Jordan2024} & hypothèse \\
13 & La dopamine de cellules vivantes modifie le poids d'un transistor organique de façon durable et réversible & Des cellules sécrétant de la dopamine ont été cultivées au-dessus d'un composant neuromorphique organique; l'oxydation de la dopamine sur l'électrode modifie la conductance du canal; un changement durable du poids et son retour ont été démontrés & \cite{Keene2020} & mesuré \\
14 & Des organoïdes contenant des neurones \nt{DOPA} fonctionnels existent; leur dopamine n'a pas été utilisée comme professeur & Dans des organoïdes issus de cellules souches humaines, on a trouvé des neurones \nt{DOPA} matures électriquement actifs et une production de dopamine. Nous n'avons trouvé dans la littérature aucune expérience où cette dopamine enseigne à un autre réseau & \cite{Jo2016} & mesuré \\
15 & Un organoïde tient un pendule en équilibre: les stimuli d'apprentissage choisis par le programme font mieux que le hasard, sans consolidation, via les synapses \nt{GLUT} & Organoïdes de cortex de souris en boucle fermée sur la tâche du \og pendule inversé sur chariot\fg{}: pour la plupart des organoïdes, les signaux choisis par apprentissage par renforcement ont donné un meilleur résultat que des signaux aléatoires ou l'absence de signal; après $45$~min de repos, l'amélioration ne subsistait pas; le blocage des récepteurs AMPA et NMDA la supprimait & \cite{Robbins2026} & mesuré \\
16 & Sans registres de contrôle, le réseau du banc ne peut pas décider lui-même de ce qui compte comme un succès (proposition~\ref{prop:livingmachine}) & Dans toutes les expériences des lignes 5--15, le signal d'apprentissage est fixé par l'expérimentateur ou par un programme; aucune expérience ne montre une culture qui élaborerait elle-même un critère de succès. C'est une conclusion tirée d'une absence, et non une expérience & --- & hypothèse \\
\bottomrule
\end{longtable}}

\section{Chapitre~\ref{sec:dishbrain}. DishBrain}

La conception du banc et les résultats du jeu proviennent d'une seule étude expérimentale et de sa suite; les formules du bruit et de la dérive vers les bons réglages sont établies dans le chapitre et vérifiées sur le modèle du livre, tandis que le mécanisme d'apprentissage dans la boîte n'est pas établi.

{\footnotesize
\setlength{\tabcolsep}{3pt}\renewcommand{\arraystretch}{1.15}
\begin{longtable}{>{\raggedright}p{0.5cm}>{\raggedright}p{4.0cm}>{\raggedright}p{5.6cm}>{\raggedright}p{2.3cm}>{\raggedright\arraybackslash}p{1.7cm}}
\toprule
N\textsuperscript{o} & Affirmation & Expérience et ce qui est mesuré & Source & Statut \\
\midrule\endfirsthead
\toprule
N\textsuperscript{o} & Affirmation & Expérience et ce qui est mesuré & Source & Statut \\
\midrule\endhead
1 & Banc: environ $800\,000$ cellules de cortex de souris ou environ $10^6$ cellules humaines par puce; $26\,000$ électrodes sur $8$~mm$^2$, enregistrement jusqu'à $1024$, stimulation par $8$ & Méthodes de l'étude: puce MaxOne, $26\,000$ électrodes de platine sur $8$~mm$^2$, enregistrement de $1024$ canaux à $20\,000$ échantillons par seconde; techniquement, on peut stimuler jusqu'à $32$ électrodes, $8$ ont pu l'être indépendamment; les cultures de souris ont été utilisées à partir d'environ $10$~jours & \cite{Kagan2022} & mesuré \\
2 & Boucle cadencée à $10$~ms: les impulsions des zones motrices déplacent la raquette (fig.~\ref{fig:dishloop}) & Les impulsions sont comptées par fenêtres de $10$~ms; le délai entre une impulsion et le stimulus de réponse, d'environ $5$~ms, est fixé par le tampon du matériel & \cite{Kagan2022} & mesuré \\
3 & Entrée: code de place (laquelle des $8$ électrodes) et fréquence de $4$ à $40$~imp./s & Dans l'expérience principale, la fréquence de stimulation croît linéairement de $4$~Hz quand la balle est près du mur du fond à $40$~Hz près de la raquette; l'amplitude des impulsions de la balle est de $75$~mV; les impulsions sont rectangulaires biphasiques & \cite{Kagan2022} & mesuré \\
4 & Sortie $\Delta p=c\,(n_\uparrow-n_\downarrow)$~\eqref{eq:dishreadout}, bruit du compte par fenêtre $1/\sqrt{RT}$ & La règle de différence entre deux zones est décrite dans l'étude (avec des poids des zones selon leur activité), la formule exacte n'est pas donnée. L'estimation du bruit découle du comptage poissonien et est établie dans le chapitre & \cite{Kagan2022}; démonstration dans le chapitre & théorie \\
5 & Professeur: après un renvoi, stimulus prévisible de $75$~mV, $100$~imp./s, $0.1$~s; après un raté, stimulus imprévisible de $150$~mV, $5$~imp./s, $4$~s en des lieux et instants aléatoires, puis $4$~s de pause & Méthodes: après un renvoi, $75$~mV, $100$~Hz, $100$~ms sur les $8$ électrodes à la fois; après un raté, $150$~mV, $5$~Hz sur des électrodes aléatoires à des instants aléatoires pendant $4$~s, puis $4$~s sans stimulation. La culture ne contient pas de dopamine & \cite{Kagan2022} & mesuré \\
6 & Le réseau agit de façon à rendre son entrée prévisible & Le principe de l'énergie libre est une proposition théorique dont les auteurs ont tiré le protocole; l'expérience est compatible avec lui, mais ne le distingue pas de mécanismes plus simples (lignes~7--8) & \cite{Friston2010,Kagan2022} & hypothèse \\
7 & Si l'échec secoue le réseau et non le succès, le temps passé dans un réglage vaut $\rho\propto1/P_{\text{raté}}$~\eqref{eq:dishdensity} (proposition~\ref{prop:dishscramble}) & Découle de l'équilibre des flux dans une chaîne de Markov à secousses symétriques, établi dans le chapitre. Aucun test direct sur une culture & démonstration dans le chapitre & théorie \\
8 & Le modèle du \og brassage après un raté\fg{} apprend: proportion de renvois $0.21\to0.38$; avec une secousse après chaque échange $\approx0.12$, sans secousse $0.19$ (fig.~\ref{fig:dishmodel}) & Calcul, script \texttt{models/dishbrain\_model.py}, $40$ cultures modèles de $2000$ échanges chacune: $0.21\to0.38$; $0.13\to0.11$; $0.19\to0.19$ & --- & modèle \\
9 & Les séries de renvois ne s'allongent que chez les cultures vivantes en boucle complète; les contrôles n'apprennent pas & $399$ sessions de $20$~min: $9$ cultures de souris ($101$ sessions), $11$ humaines ($138$), $6$ puces avec milieu seul ($80$), $20$ cultures au repos ($42$), $3$ simulations ($38$). La longueur moyenne d'un échange sur les $15$ dernières minutes n'est supérieure à celle des $5$ premières que chez les cultures vivantes; les cellules non neuronales (HEK293T) et les puces avec milieu seul ne montrent aucun effet & \cite{Kagan2022} & mesuré \\
10 & Hausse significative en une session seulement avec rétroaction complète; avec du silence à la place du stimulus, le jeu est meilleur en fin de session que sans rétroaction, mais l'effet est moindre & $486$ sessions sur $3$ jours: \og stimulus\fg{} ($n=27$), \og silence\fg{} ($17$), \og sans rétroaction\fg{} ($15$). La hausse au fil du temps n'est significative que dans la condition \og stimulus\fg{}; en fin de session, la condition \og silence\fg{} dépassait \og sans rétroaction\fg{} avec un effet moindre & \cite{Kagan2022} & mesuré \\
11 & L'effet est faible; savoir si le réseau apprend durablement reste une question ouverte & Au sein d'une session, l'amélioration est robuste, mais, de l'aveu même des auteurs, aucun apprentissage stable d'une session à l'autre n'a été observé sur plusieurs jours & \cite{Kagan2022} & mesuré \\
12 & Pendant le jeu, le réseau se rapproche du régime critique, et plus il s'en approche, meilleur est le jeu & Analyse des avalanches d'activité des mêmes cultures: une entrée structurée rapproche le réseau du régime critique, et la qualité du jeu est corrélée à cette proximité; mais la criticité seule, sans information sur les conséquences des actions, ne suffit pas à l'apprentissage & \cite{Habibollahi2023} & mesuré \\
13 & La \og sensibilité\fg{} de la culture n'est pas démontrée & Article de réponse de trente chercheurs: un changement de comportement en boucle fermée ne prouve ni intelligence ni sensations; aucune expérience ne le démontre & \cite{Balci2023} & hypothèse \\
14 & Quatrième contrôle proposé: un stimulus prévisible après un raté, de même durée et de même énergie (section~\ref{sec:dishspec}) & Cette expérience n'a pas été réalisée; c'est une proposition du livre & --- & hypothèse \\
\bottomrule
\end{longtable}}

